%Version 3.1 December 2024
% See section 11 of the User Manual for version history
%
%%%%%%%%%%%%%%%%%%%%%%%%%%%%%%%%%%%%%%%%%%%%%%%%%%%%%%%%%%%%%%%%%%%%%%
%%                                                                 %%
%% Please do not use \input{...} to include other tex files.       %%
%% Submit your LaTeX manuscript as one .tex document.              %%
%%                                                                 %%
%% All additional figures and files should be attached             %%
%% separately and not embedded in the \TeX\ document itself.       %%
%%                                                                 %%
%%%%%%%%%%%%%%%%%%%%%%%%%%%%%%%%%%%%%%%%%%%%%%%%%%%%%%%%%%%%%%%%%%%%%

%%\documentclass[referee,sn-basic]{sn-jnl}% referee option is meant for double line spacing

%%=======================================================%%
%% to print line numbers in the margin use lineno option %%
%%=======================================================%%

%%\documentclass[lineno,pdflatex,sn-basic]{sn-jnl}% Basic Springer Nature Reference Style/Chemistry Reference Style

%%=========================================================================================%%
%% the documentclass is set to pdflatex as default. You can delete it if not appropriate.  %%
%%=========================================================================================%%

%%\documentclass[sn-basic]{sn-jnl}% Basic Springer Nature Reference Style/Chemistry Reference Style

%%Note: the following reference styles support Namedate and Numbered referencing. By default the style follows the most common style. To switch between the options you can add or remove “Numbered” in the optional parenthesis. 
%%The option is available for: sn-basic.bst, sn-chicago.bst%  
 
%%\documentclass[pdflatex,sn-nature]{sn-jnl}% Style for submissions to Nature Portfolio journals
%%\documentclass[pdflatex,sn-basic]{sn-jnl}% Basic Springer Nature Reference Style/Chemistry Reference Style
\documentclass[pdflatex,sn-mathphys-num]{sn-jnl}% Math and Physical Sciences Numbered Reference Style
%%\documentclass[pdflatex,sn-mathphys-ay]{sn-jnl}% Math and Physical Sciences Author Year Reference Style
%%\documentclass[pdflatex,sn-aps]{sn-jnl}% American Physical Society (APS) Reference Style
%%\documentclass[pdflatex,sn-vancouver-num]{sn-jnl}% Vancouver Numbered Reference Style
%%\documentclass[pdflatex,sn-vancouver-ay]{sn-jnl}% Vancouver Author Year Reference Style
%%\documentclass[pdflatex,sn-apa]{sn-jnl}% APA Reference Style
%%\documentclass[pdflatex,sn-chicago]{sn-jnl}% Chicago-based Humanities Reference Style

%%%% Standard Packages
%%<additional latex packages if required can be included here>

\usepackage{graphicx}%
\usepackage{multirow}%
\usepackage{amsmath,amssymb,amsfonts}%
\usepackage{amsthm}%
\usepackage{mathrsfs}%
\usepackage[title]{appendix}%
\usepackage{xcolor}%
\usepackage{textcomp}%
\usepackage{manyfoot}%
\usepackage{booktabs}%
\usepackage{algorithm}%
\usepackage{algorithmicx}%
\usepackage{algpseudocode}%
\usepackage{listings}%
\usepackage{placeins}%
%%%%

\usepackage{mathtools} 

% Commands carried over from the former INFORMS manuscript
\newcommand{\re}{\mathbb{R}}
\newcommand{\N}{\mathbb{N}}
\newcommand{\lmd}{\lambda}
\DeclareMathOperator{\rank}{rank}
\newcommand{\reff}[1]{(\ref{#1})}
\newcommand{\rd}[1]{{\color{red}#1}} 
\newcommand{\bl}[1]{{\color{blue}#1}}
\newcommand{\rdurl}[1]{\href{#1}{\textcolor{red}{\nolinkurl{#1}}}}
\newcommand{\vect}[1]{\boldsymbol{#1}}
\newcommand{\mat}[1]{\boldsymbol{#1}}
\newcommand{\supp}[1]{\operatorname{supp}(#1)}
\newcommand{\st}{\mathrm{s.t.}}
\newcommand{\ddd}{,\ldots,}

%%%%%=============================================================================%%%%
%%%%  Remarks: This template is provided to aid authors with the preparation
%%%%  of original research articles intended for submission to journals published 
%%%%  by Springer Nature. The guidance has been prepared in partnership with 
%%%%  production teams to conform to Springer Nature technical requirements. 
%%%%  Editorial and presentation requirements differ among journal portfolios and 
%%%%  research disciplines. You may find sections in this template are irrelevant 
%%%%  to your work and are empowered to omit any such section if allowed by the 
%%%%  journal you intend to submit to. The submission guidelines and policies 
%%%%  of the journal take precedence. A detailed User Manual is available in the 
%%%%  template package for technical guidance.
%%%%%=============================================================================%%%%

%% as per the requirement new theorem styles can be included as shown below
\theoremstyle{thmstyleone}%
\newtheorem{theorem}{Theorem}%  meant for continuous numbers
%%\newtheorem{theorem}{Theorem}[section]% meant for sectionwise numbers
%% optional argument [theorem] produces theorem numbering sequence instead of independent numbers for Proposition
\newtheorem{proposition}[theorem]{Proposition}% 
\newtheorem{corollary}[theorem]{Corollary}
\newtheorem{assumption}{Assumption}
%%\newtheorem{proposition}{Proposition}% to get separate numbers for theorem and proposition etc.

\theoremstyle{thmstyletwo}%
\newtheorem{example}{Example}%
\newtheorem{remark}{Remark}%

\theoremstyle{thmstylethree}%
\newtheorem{definition}{Definition}%

\raggedbottom
%%\unnumbered% uncomment this for unnumbered level heads

\begin{document}

\title[Log-Polynomial Optimization]{Log-Polynomial Optimization}

%%=============================================================%%
%% GivenName	-> \fnm{Joergen W.}
%% Particle	-> \spfx{van der} -> surname prefix
%% FamilyName	-> \sur{Ploeg}
%% Suffix	-> \sfx{IV}
%% \author*[1,2]{\fnm{Joergen W.} \spfx{van der} \sur{Ploeg} 
%%  \sfx{IV}}\email{iauthor@gmail.com}
%%=============================================================%%

\author[1]{\fnm{Jiyoung} \sur{Choi}}
\email{jichoi@ksu.edu}
\equalcont{These authors contributed equally to this work.}

\author[2]{\fnm{Jiawang} \sur{Nie}}
\email{njw@math.ucsd.edu}
\equalcont{These authors contributed equally to this work.}

\author[3]{\fnm{Xindong} \sur{Tang}}
\email{xdtang@hkbu.edu.hk}
\equalcont{These authors contributed equally to this work.}

\author*[4]{\fnm{Suhan} \sur{Zhong}}
\email{suzhong@sjtu.edu.cn}
\equalcont{These authors contributed equally to this work.}

\affil[1]{\orgdiv{Department of Mathematics},
    \orgname{Kansas State University},
    \orgaddress{\city{Manhattan}, \state{KS} \postcode{66506}, \country{USA}}}

\affil[2]{\orgdiv{Department of Mathematics},
    \orgname{University of California San Diego},
\orgaddress{\street{9500 Gilman Drive}, \city{La Jolla}, \state{CA} \postcode{92093}, \country{USA}}}

\affil[3]{\orgdiv{Department of Mathematics},
    \orgname{Hong Kong Baptist University},
    \orgaddress{\city{Kowloon Tong}, \country{Hong Kong}}}

\affil*[4]{\orgdiv{School of Mathematical Sciences},
    \orgname{Shanghai Jiao Tong University},
    \orgaddress{\city{Shanghai}, \country{China}}}

%%==================================%%
%% Sample for unstructured abstract %%
%%==================================%%

\abstract{We study a class of optimization problems in which the objective is
given as a \rd{weighted sum of logarithms of polynomial functions}.
Motivating applications include maximum likelihood estimation,
cross-entropy, and Kullback--Leibler divergence
when the target distribution is fixed and the predicted
probabilities are polynomial in the decision variables.
We propose a hierarchy of moment relaxations based on
truncated $K$-moment problems for log-polynomial optimization.
\rd{By preserving the logarithmic structure, the relaxation order is
independent of the positive weights.}
We distinguish the convexification gap from the
moment-truncation gap and establish asymptotic convergence to the
convexified moment relaxation under the Archimedean condition.
\rd{We provide sufficient conditions for tightness of the hierarchy.
We also give conditions under which global optimizers can be extracted
from optimal moment sequences.}
We further strengthen the relaxations using
Lagrange multiplier expressions.
Applications and numerical experiments
quantify both the practical scope and the limitations of the method
\rd{across representative problem classes and parameter regimes}.}

%%================================%%
%% Sample for structured abstract %%
%%================================%%

% \abstract{\textbf{Purpose:} The abstract serves both as a general introduction to the topic and as a brief, non-technical summary of the main results and their implications. The abstract must not include subheadings (unless expressly permitted in the journal's Instructions to Authors), equations or citations. As a guide the abstract should not exceed 200 words. Most journals do not set a hard limit however authors are advised to check the author instructions for the journal they are submitting to.
% 
% \textbf{Methods:} The abstract serves both as a general introduction to the topic and as a brief, non-technical summary of the main results and their implications. The abstract must not include subheadings (unless expressly permitted in the journal's Instructions to Authors), equations or citations. As a guide the abstract should not exceed 200 words. Most journals do not set a hard limit however authors are advised to check the author instructions for the journal they are submitting to.
% 
% \textbf{Results:} The abstract serves both as a general introduction to the topic and as a brief, non-technical summary of the main results and their implications. The abstract must not include subheadings (unless expressly permitted in the journal's Instructions to Authors), equations or citations. As a guide the abstract should not exceed 200 words. Most journals do not set a hard limit however authors are advised to check the author instructions for the journal they are submitting to.
% 
% \textbf{Conclusion:} The abstract serves both as a general introduction to the topic and as a brief, non-technical summary of the main results and their implications. The abstract must not include subheadings (unless expressly permitted in the journal's Instructions to Authors), equations or citations. As a guide the abstract should not exceed 200 words. Most journals do not set a hard limit however authors are advised to check the author instructions for the journal they are submitting to.}

\keywords{Log-polynomial optimization, Polynomial optimization,
Truncated $K$-moment problem, Moment relaxations,
Semidefinite programming, Lagrange multiplier expressions}

%%\pacs[JEL Classification]{D8, H51}

%%\pacs[MSC Classification]{35A01, 65L10, 65L12, 65L20, 65L70}

\maketitle
\section{Introduction}\label{sec:Intro}


Consider the following log-polynomial optimization problem:
\begin{equation}\label{eq:problem}
	\left\{\begin{array}{cl}
		\max\limits_{\vect{x}\in\re^n} &
		f(\vect{x}) = \sum\limits_{i=1}^m a_i\log p_i(\vect{x})  \\
		\st & c_j(\vect{x}) = 0 \,\, (j \in \mathcal{E}),\\
		& c_j(\vect{x}) \geq 0 \,\, (j \in \mathcal{I}),
	\end{array}\right.
\end{equation}
where $a_1,\ldots, a_m$ are positive scalars, 
$\mathcal{E}$ and $\mathcal{I}$ are two finite disjoint index sets
and $p_1,\ldots,p_m$ are polynomials, while each $c_j$ is a polynomial for $j\in \mathcal{E}\cup\mathcal{I}$.
Denote $\vect{c}_{\mathcal{E}} = (c_j)_{j\in\mathcal{E}}$ and $\vect{c}_{\mathcal{I}} = (c_j)_{j\in\mathcal{I}}$.
The feasible set of \eqref{eq:problem} can be written as
\begin{equation}\label{eq:K}
	K  \coloneqq  \{ \vect{x}\in\re^n : \vect{c}_{\mathcal{E}}(\vect{x}) = 0,\, \vect{c}_{\mathcal{I}}(\vect{x}) \ge  0\}.
\end{equation}
\begin{assumption}\label{ass:positivity}
\rd{For every $i\in[m]$, $p_i$ is nonnegative on $K$. Moreover, there
exists $\vect{x}^{\circ}\in K$ such that $p_i(\vect{x}^{\circ})>0$ for every $i\in[m]$.}
\end{assumption}
\rd{We impose Assumption~\ref{ass:positivity} throughout this paper.}
We use the extended-value convention $\log 0=-\infty$.
For convenience, denote
\begin{equation}\label{eq:weighted_product}
p^{\vect{a}}(\vect{x}) \coloneqq p_1(\vect{x})^{a_1}\cdots p_m(\vect{x})^{a_m}
= \exp(f(\vect{x})).
\end{equation}
The log-polynomial optimization \eqref{eq:problem}
is typically nonconvex.
If all $a_i$ are positive integers, 
then $p^{\vect{a}}$ is a polynomial.
In this case, the problem \eqref{eq:problem} 
is equivalent to maximizing $p^{\vect{a}}$ over $K$.
The resulting polynomial optimization problem can be addressed
using the standard moment--sum-of-squares (Moment--SOS) hierarchy
\cite{lasserre2001,MPO}. However, its computational cost can increase
significantly when the integer exponents are large. If some $a_i$ are
noninteger, $p^{\vect{a}}$ is generally not a polynomial, so this
direct polynomial-product reformulation is unavailable.
Consequently, when the coefficients are general positive scalars,
it can be difficult to compute a useful upper bound
or certify global optimality for \eqref{eq:problem}.



Log-polynomials have many applications.
In combinatorial optimization, 
log-concave polynomials serve as an important tool 
to study counting problems \cite{Anari24}. 
In statistical learning, many common objectives and loss functions,
such as log-likelihood and cross-entropy,
are defined with logarithmic terms.
These functions are widely used in density estimation
\cite{Carpenter2018,Dempster1977,Fisher1922,Vinayak2019,Silverman1986}
and classification tasks
\cite{Bishop2006,Goodfellow2016,Goodfellow2014,Mannor2005,Mao2023,Zhang2018},
where learning often incorporates structural requirements such as non-negativity and normalization constraints, physics-informed or conservation-law constraints \cite{Gusmao2024,Hansen2024,Rayas2022},
domain-knowledge constraints \cite{Karanam2025}, total-positivity restrictions \cite{Lauritzen2019}, or sparsity-inducing regularization
\cite{Banerjee2008,Srivastava2024}.
Suppose \rd{the selected model probabilities are polynomial functions
of $\vect{x}$, the coefficients multiplying the logarithms are fixed
scalars, and the structural constraints define a semialgebraic feasible
set $K$}.
Then this subclass of the learning problem reduces to solving
log-polynomial optimization.
We illustrate this framework via the following motivating applications.


{\it Maximum likelihood estimation (MLE)} is a core problem in statistics.
\rd{Consider a data set $\mathcal{D}$ generated from a polynomial
probability model $p(\xi; \vect{x})$, where $p(\xi;\vect{x})$ is
polynomial in $\vect{x}$ and the feasible set $K$ restricts
$\vect{x}$ so that $p(\cdot;\vect{x})$ is a valid probability mass
function.}
\rd{Let $S = \{\xi^{(1)}, \ldots, \xi^{(m)}\}$
denote the set of distinct observations in $\mathcal{D}$.}
We use $a_i$ to count the number of samples in $\mathcal{D}$ that equal $\xi^{(i)}$
for $i=1,\ldots, m$.
The {\em likelihood function} associated with $\mathcal{D}$ 
and the corresponding {\em log-likelihood function} are respectively 
defined as
\[
L(\mathcal{D};\vect{x}) \coloneqq \prod\limits_{i=1}^m p(\xi^{(i)};\vect{x})^{a_i},
\quad 
\log L(\mathcal{D}; \vect{x}) = \sum\limits_{i=1}^m a_i\log p(\xi^{(i)};\vect{x}).
\]
The MLE aims to find parameter values that maximize the likelihood function,
which is equivalent to maximizing the log-likelihood.

{\it Cross-entropy} and {\it Kullback--Leibler (KL) divergence}
are commonly used in statistical learning and probabilistic modeling.
Consider the discrete random variable $\xi$ supported on 
$S = \{\xi^{(1)}, \ldots, \xi^{(m)}\}$, 
whose true distribution $P$ has $p(\xi)$ as its {\em probability mass function} (PMF).
Let $Q(\vect{x})$ be the predicted distribution with a polynomial PMF $q(\xi; \vect{x})$.
The {\em cross-entropy} between $P$ and $Q(\vect{x})$ is
\[
H(P,Q(\vect{x})) \coloneqq -\sum\limits_{i=1}^m p(\xi^{(i)})\log q(\xi^{(i)}; \vect{x}).
\]
The KL divergence from $P$ to $Q(\vect{x})$ is 
\[
D_{\operatorname{KL}}(P\,\|\,Q(\vect{x})) \coloneqq 
\sum\limits_{i=1}^m p(\xi^{(i)})
\big(\log p(\xi^{(i)})-\log q(\xi^{(i)};\vect{x})\big).
\] 
In a learning process, the goal is to determine an optimal parameter
$\vect{x}\in K$ such that $H(P, Q(\vect{x}))$ or $D_{\operatorname{KL}}(P\,\|\,Q(\vect{x}))$
is minimized.
Since $P$ is fixed and $S$ is its support, the values
$p(\xi^{(i)})$ are positive fixed scalars. Hence, minimizing either
cross-entropy or KL divergence is equivalent to \eqref{eq:problem}
with $a_i=p(\xi^{(i)})$ and
$p_i(\vect{x})=q(\xi^{(i)};\vect{x})$.
By contrast, the entropy
$-\sum_i q(\xi^{(i)};\vect{x})\log q(\xi^{(i)};\vect{x})$ has
decision-dependent coefficients and therefore lies outside the scope of
\eqref{eq:problem}.

\rd{Bach~\cite{Bach2025} develops sum-of-squares relaxations
for estimating $f$-divergences from moment
information, computing normalizing integrals, and performing
variational inference.
By contrast, the present work studies the global optimization of
logarithms of polynomial functions with fixed scalar weights over
semialgebraic sets. Our hierarchy yields upper bounds on the global
optimum; we also provide sufficient conditions for finite-order
exactness, procedures for extracting global optimizers under
appropriate flat-truncation conditions, and strengthened relaxations
based on Lagrange multiplier expressions.}



\subsection{Contributions}
In this paper, we propose a hierarchy of moment relaxations 
for log-polynomial optimization.
\rd{When all $a_i$ are positive integers, one can apply the standard
Moment--SOS hierarchy to the polynomial $p^{\vect{a}}$ defined in
\eqref{eq:weighted_product}. However, the degree of this polynomial
grows with the integer exponents, and this direct polynomial-product
reformulation is generally unavailable when some $a_i$ are
noninteger.}
\rd{Our formulation retains the individual polynomials $p_i$ and has
a relaxation order that does not depend on the weights $a_i$.}
\rd{We construct a moment relaxation hierarchy for
\eqref{eq:problem}, where each relaxation problem has an objective
consisting of logarithms of linear functions of the moment variables
and has semidefinite constraints.}
\rd{We distinguish the convexification gap
$f_{\operatorname{mom}}-f_{\max}$ from the moment-truncation gap $f_{\operatorname{mom},k}-f_{\operatorname{mom}}$.}
\rd{We study the properties of these relaxations and provide
sufficient conditions under which
$f_{\operatorname{mom},k}=f_{\max}$ at a finite relaxation order.}
\rd{Under an appropriate flat-truncation condition, we provide
procedures for extracting global optimizers.}
\rd{In particular, we analyze a class in which the objective
polynomials and inequality-constraint polynomials are SOS-concave and
the equality constraints are linear.}



This paper is organized as follows.
Section~\ref{sec:Pre} introduces the notation and some preliminaries.
Section~\ref{sec:rel} presents the proposed moment relaxation 
approach based on the truncated $K$-moment problems.
In Section~\ref{sec:tight}, we provide theoretical guarantees for 
the relaxation's tightness and discuss conditions under 
which the global optimizers can be extracted.
In Section~\ref{sec:LME}, we introduce Lagrange multiplier expressions (LMEs)
to obtain strengthened moment relaxations.
Numerical experiments and applications are presented in Section~\ref{sec:exp}.
Section~\ref{sec:con} concludes the paper with a discussion of the
results and limitations.



\section{Preliminaries}\label{sec:Pre}


\subsection{Notation}
The symbol $\mathbb{N}$ (resp., $\mathbb{R}$) represents the set of 
nonnegative integers (resp., real numbers).
For a positive integer $n$, $[n]\coloneqq \{1,\ldots, n\}$ and 
$\mathbb{R}^n$ denotes the $n$-dimensional real Euclidean space.

For a set $S\subseteq\mathbb{R}^n$, $\operatorname{conv}(S)$ denotes
the convex hull of $S$.

For a vector $\vect{u} \in \mathbb{R}^n$, $\|\vect{u}\|$ denotes the standard 
Euclidean norm and $\delta_{\vect{u}}$ denotes the Dirac measure supported at $\vect{u}$.
The vectors $\mathbf{1}_n,\mathbf{0}_n\in\mathbb{R}^n$
denote the all-ones and all-zeros vectors, respectively.
Their subscripts may be omitted when the dimension is clear.
For a real symmetric matrix $\mat{A}$, the notation $\mat{A}\succeq0$
means that $\mat{A}$ is positive semidefinite (PSD).
Let $\vect{x} \coloneqq (x_1, \ldots, x_n)$ be a vector of variables.
We use $\mathbb{R}[\vect{x}]$ to denote the ring of polynomials 
with real coefficients in $\vect{x}$,
and $\mathbb{R}[\vect{x}]_d$ its subset of polynomials
of degree at most $d$.
For a function $f(\vect{x})$, $\nabla f(\vect{x})$ denotes its full gradient
and $\frac{\partial f}{\partial x_i}(\vect{x})$ denotes its partial
derivative with respect to $x_i$.
For $\vect{\alpha} \coloneqq (\alpha_1, \ldots, \alpha_n)\in \mathbb{N}^n$,
denote $\vect{x}^{\vect{\alpha}} \coloneqq x_1^{\alpha_1} \cdots x_n^{\alpha_n}$,
whose degree is counted by $|\vect{\alpha}| \coloneqq \alpha_1 + \cdots + \alpha_n$.
Given a degree $d$, we denote the multi-index set
\[
\mathbb{N}_d^n \coloneqq \{\vect{\alpha} \in \mathbb{N}^n : |\vect{\alpha}| \leq d \} .
\]
The column vector of all monomials in $\vect{x}$ with degrees up to $d$ is denoted as
\begin{equation}\label{eq:xvec}
	[\vect{x}]_d \coloneqq \begin{bmatrix}
		1 & x_1 & \ldots & x_n & x_1^2 & x_1x_2 & \ldots & x_n^d
	\end{bmatrix}^T.
\end{equation}



\subsection{Nonnegative polynomials}

For a polynomial $p\in\mathbb{R}[\vect{x}]$ and subsets
$I,J\subseteq\mathbb{R}[\vect{x}]$, define
\begin{equation*}
	p \cdot I \coloneqq \{pq : q \in I\} 
	\quad \text{and} \quad 
	I+J \coloneqq \{a+b : a \in I, b \in J\}.
\end{equation*}
A nonempty subset $I \subseteq \mathbb{R}[\vect{x}]$ is an \textit{ideal}
of $\mathbb{R}[\vect{x}]$ if $I+I \subseteq I$ and $\mathbb{R}[\vect{x}] \cdot I \subseteq I$.
A polynomial $\sigma \in \mathbb{R}[\vect{x}]$ is a \textit{sum of squares (SOS)} 
if $\sigma = q_1^2 + \cdots + q_k^2$ for some polynomials $q_i \in \mathbb{R}[\vect{x}]$. 
The set of all SOS polynomials is denoted by $\Sigma[\vect{x}]$.

Let $\vect{h} = (h_1, \ldots, h_{\nu})$ and
$\vect{g}=(g_1,\ldots,g_{\rho})$
be two polynomial tuples. 
The real zero set of $\vect{h}$ is the set 
$Z(\vect{h}) \coloneqq \{\vect{x} \in \mathbb{R}^n:
h_1(\vect{x}) = \cdots = h_{\nu}(\vect{x}) = 0\}$.
The ideal generated by $\vect{h}$ is
\begin{equation*}
	\operatorname{Ideal}[\vect{h}]\, \coloneqq \,
	h_1 \cdot \mathbb{R}[\vect{x}] + \cdots
	+ h_{\nu} \cdot \mathbb{R}[\vect{x}] .
\end{equation*}
It is clear that every $p\in
\operatorname{Ideal}[\vect{h}]$ vanishes on $Z(\vect{h})$.
Consider the semialgebraic set determined by $\vect{g}$:
\begin{equation*}
	S(\vect{g}) \coloneqq \{\vect{x} \in \mathbb{R}^n :
	g_1(\vect{x}) \geq 0, \ldots,
	g_{\rho}(\vect{x}) \geq 0\}.
\end{equation*}
If a polynomial $f$ admits a decomposition
\begin{equation}\label{eq:fsos}
	f = \sigma_0 + \sigma_1 g_1 + \cdots
	+ \sigma_{\rho}g_{\rho},
\end{equation}
for $\sigma_0,\ldots,\sigma_{\rho}\in\Sigma[\vect{x}]$,
then $f$ is non-negative on $S(\vect{g})$.
The \textit{quadratic module} generated by $\vect{g}$ consists of all polynomials
in form of \eqref{eq:fsos}, which reads
\begin{equation*}
	\operatorname{QM}[\vect{g}] \coloneqq
	\Sigma[\vect{x}] + g_1 \cdot \Sigma[\vect{x}]
	+ \cdots + g_{\rho} \cdot \Sigma[\vect{x}].
\end{equation*}
If $p\in \operatorname{Ideal}[\vect{h}] + \operatorname{QM}[\vect{g}]$,
then $p(\vect{x})\ge 0$ for all $\vect{x}\in Z(\vect{h}) \cap S(\vect{g})$.
But the converse is not always true.
The set $\operatorname{Ideal}[\vect{h}] + \operatorname{QM}[\vect{g}]$ 
is said to be \textit{Archimedean} if it contains a polynomial $q$ 
such that $\{\vect{x} \in \mathbb{R}^n : q(\vect{x}) \ge 0\}$ is a compact set.
Under this condition, if a polynomial $f$ is strictly 
positive on $Z(\vect{h})\cap S(\vect{g})$, 
then it must belong to $\operatorname{Ideal}[\vect{h}]+\operatorname{QM}[\vect{g}]$.
This result is called \textit{Putinar's Positivstellensatz} \cite{Putinar}.


\subsection{Moment and localizing matrices}
\label{ssc:momloc}

The space of real vectors indexed by $\vect{\alpha} \in \mathbb{N}_{2k}^n$ 
is denoted by $\mathbb{R}^{\mathbb{N}_{2k}^n}$. 
A vector $\vect{y} = (y_{\vect{\alpha}})_{\vect{\alpha} \in \mathbb{N}_{2k}^n}$ 
in this space is called a \textit{truncated multi-sequence} of degree $2k$.
For an integer $s$ with $0\le s\le k$, 
the $s$-th order \textit{moment matrix} generated by $\vect{y}$ is defined as
\begin{equation*}
	\mat{M}_{s}[\vect{y}] \coloneqq
	[y_{\vect{\alpha}+\vect{\beta}}]_{\vect{\alpha},\vect{\beta}\in\mathbb{N}_{s}^n}.
\end{equation*}
The rows and columns of $\mat{M}_k[\vect{y}]$ are indexed by multi-indices 
from $\mathbb{N}_{k}^n$, typically arranged in graded lexicographic order.
For example, when $n = 2$ and $k = 2$, the matrix $\mat{M}_2[\vect{y}]$ is given by
\begin{equation*}
	\mat{M}_2[\vect{y}] = \begin{bmatrix}
		y_{00} & y_{10} & y_{01} & y_{20} & y_{11} & y_{02} \\
		y_{10} & y_{20} & y_{11} & y_{30} & y_{21} & y_{12} \\
		y_{01} & y_{11} & y_{02} & y_{21} & y_{12} & y_{03} \\
		y_{20} & y_{30} & y_{21} & y_{40} & y_{31} & y_{22} \\
		y_{11} & y_{21} & y_{12} & y_{31} & y_{22} & y_{13} \\
		y_{02} & y_{12} & y_{03} & y_{22} & y_{13} & y_{04}
	\end{bmatrix}.
\end{equation*}
Every polynomial $f\in\re[\vect{x}]_{2k}$ can be identified with its coefficient vector
$\operatorname{vec}(f) \coloneqq
(f_{\vect{\alpha}})_{\vect{\alpha}\in \mathbb{N}_{2k}^n}$,
i.e., $f(\vect{x}) =  \operatorname{vec}(f)^T[\vect{x}]_{2k}$. 
Thus, a truncated multi-sequence $\vect{y} \in \mathbb{R}^{\mathbb{N}_{2k}^n}$ defines a 
linear functional on $\mathbb{R}[\vect{x}]_{2k}$ as
\begin{equation}\label{eq:bilinear}
	\langle f, \vect{y} \rangle \coloneqq 
	\sum_{\vect{\alpha} \in \mathbb{N}_{2k}^n}
	f_{\vect{\alpha}}y_{\vect{\alpha}}
	= \operatorname{vec}(f)^T \vect{y}.
\end{equation}
For a polynomial $q \in \mathbb{R}[\vect{x}]_{2k}$, 
let $s \coloneqq k - \lceil \deg(q)/2 \rceil$.
The product $q(\vect{x}) [\vect{x}]_s [\vect{x}]_s^T$ is a $\binom{n+s}{s}$-by-$\binom{n+s}{s}$
symmetric polynomial matrix. 
Its expansion can be written as
\begin{equation*}
	q(\vect{x}) [\vect{x}]_s [\vect{x}]_s^T
	= \sum_{\vect{\alpha} \in \mathbb{N}_{2k}^n}
	\vect{x}^{\vect{\alpha}}\mat{Q}_{\vect{\alpha}}
\end{equation*}
for some constant symmetric matrices $\mat{Q}_{\vect{\alpha}}$. 
The $k$-th order \textit{localizing matrix} associated with 
the polynomial $q$ and the truncated multi-sequence $\vect{y} \in \mathbb{R}^{\mathbb{N}_{2k}^n}$ is defined as
\begin{equation*}
	\mat{L}^{(k)}_q [\vect{y}] \coloneqq
	\sum_{\vect{\alpha} \in \mathbb{N}_{2k}^n}
	y_{\vect{\alpha}}\mat{Q}_{\vect{\alpha}}.
\end{equation*}
In the special case $q=1$ (the constant unit polynomial),
this definition recovers the moment matrix, i.e., $\mat{L}^{(k)}_1 [\vect{y}] = \mat{M}_k[\vect{y}]$.



\subsection{Truncated \texorpdfstring{$K$}{K}-moment problems}

A truncated multi-sequence
$\vect{y} = (y_{\vect{\alpha}})_{\vect{\alpha} \in \mathbb{N}^n_{2k}}$
is said to
\textit{admit} a Borel measure $\mu$ on $K$ if
\begin{equation} \label{eq:y_admits_meas}
	y_{\vect{\alpha}} = \int_K \vect{x}^{\vect{\alpha}} d \mu
	\quad \text{for all} \quad
	\vect{\alpha} \in \mathbb{N}^n_{2k}.
\end{equation}
Such a measure $\mu$ is called a \textit{representing measure} of $\vect{y}$.
The support of $\mu$, denoted by $\supp{\mu}$,
is the smallest closed set $T \subseteq \mathbb{R}^n$ 
such that $\mu(\mathbb{R}^n \setminus T) = 0$.
A measure is called \textit{finitely atomic} if its support is a finite set.
Specifically, if $\supp{\mu}$ $= \{\vect{v}_1, \ldots, \vect{v}_r\}$,
then the measure $\mu$ is called $r$-atomic and can be written as
\[
\mu = \sum_{i=1}^r \lambda_i \delta_{\vect{v}_i} \quad 
\text{for some positive weights } \lambda_1, \dots, \lambda_r > 0.
\]
Suppose that $\deg(c_j)\leq 2k$ for every
$j\in\mathcal{E}\cup\mathcal{I}$.
For a truncated multi-sequence $\vect{y} \in \mathbb{R}^{\mathbb{N}_{2k}^n}$ to admit a representing measure
$\mu$ on $K$ as in \eqref{eq:y_admits_meas}, 
it is necessary that the following matrix conditions hold (see \cite{MPO}):
\begin{equation} \label{eq:loc_mom}
	\mat{M}_k[\vect{y}] \succeq 0, \quad \mat{L}_{c_j}^{(k)}[\vect{y}] \succeq 0 \,(j \in \mathcal{I}), \quad
	\mat{L}_{c_j}^{(k)}[\vect{y}] = 0 \, (j \in \mathcal{E}).
\end{equation}



\section{A semidefinite relaxation hierarchy}
\label{sec:rel}



In this section, we propose a moment-relaxation approach for
solving the log-polynomial optimization problem
\eqref{eq:problem}.

\subsection{A convex relaxation}
To begin with, we introduce a natural convex relaxation of the 
log-polynomial optimization \eqref{eq:problem}.
Define the degree
\begin{equation}\label{eq:d0d}
	d  \coloneqq  \max\,\Big\{ 
	\max_{i\in[m]}\left\lceil \frac{\deg(p_i)}{2}\right\rceil, 
	\, \max_{j\in\mathcal{E}}\left\lceil\frac{\deg(c_j)}{2}\right\rceil,
	\, \max_{j\in\mathcal{I}}\left\lceil\frac{\deg(c_j)}{2}\right\rceil \Big\}.
\end{equation}
Each polynomial in \eqref{eq:problem} belongs to the polynomial ring $\re[\vect{x}]_{2k}$ for all $k\ge d$.
By replacing each monomial $\vect{x}^{\vect{\alpha}}$ appearing
in the problem with an auxiliary moment variable
$y_{\vect{\alpha}}$,
we obtain the following equivalent moment reformulation of \eqref{eq:problem}:
\begin{equation}\label{eq:momequiv}
	\left\{\begin{array}{rl}
		f_{\max}\coloneqq \max\limits_{\vect{z}} & \sum\limits_{i=1}^m a_i \log \langle p_i, \vect{z}\rangle\\
		\st & \vect{z} = [\vect{x}]_{2d}\mbox{ for some $\vect{x}\in K$}.
	\end{array}
	\right.
\end{equation}
Here $\langle \cdot , \cdot \rangle$ is the bilinear operation 
defined in \eqref{eq:bilinear}, 
$[\vect{x}]_{2d}$ is the vector of monomials of degree at most
$2d$ as in \eqref{eq:xvec},
and $K$ is the feasible set given in \eqref{eq:K}.
\rd{To clarify the geometry of this reformulation, define the polynomial map
\[
\mathcal{P}:\mathbb{R}^n\to\mathbb{R}^m,\qquad
\mathcal{P}(\vect{x})\coloneqq(p_1(\vect{x}),\ldots,p_m(\vect{x})),
\]
and its linear lift to the moment space
\[
\widetilde{\mathcal{P}}:
\mathbb{R}^{\mathbb{N}_{2d}^n}\to\mathbb{R}^m,\qquad
\widetilde{\mathcal{P}}(\vect{z})
\coloneqq
\big(\langle p_1,\vect{z}\rangle,\ldots,\langle p_m,\vect{z}\rangle\big).
\]
Also, define
\[
\phi(\vect{u})\coloneqq\sum_{i=1}^m a_i\log u_i,
\qquad u_i>0\quad(i\in[m]).
\]
Then the objective function in \eqref{eq:momequiv} is
$\phi(\widetilde{\mathcal{P}}(\vect{z}))$.
Although $\phi$ is strictly concave on its domain,
$\phi\circ\widetilde{\mathcal{P}}$ need not be strictly concave
on the lifted feasible set, since distinct moment vectors may have
the same image under $\widetilde{\mathcal{P}}$.}
To obtain a convex feasible set in the lifted moment space, we
define the conic hull generated by the feasible set of
\eqref{eq:momequiv}:
\begin{equation}  \label{Rd(K)}
	\mathscr{R}_d(K) \coloneqq \Big\{ \vect{z} = \sum\limits_{j=1}^l \lambda_j [\vect{v}_j]_{2d} :
	\lambda_j \ge  0, \,\, \vect{v}_j \in K,\,\, l \in \mathbb{N} \Big\}.
\end{equation}
It follows that $\mathscr{R}_d(K)\cap \{z_0=1\}$ is the convex hull 
of the set $\{[\vect{x}]_{2d}: \vect{x}\in K\}$.
{\color{red}
The image of $K$ under $\mathcal{P}$ is
\[
\mathcal{P}(K) \coloneqq
\{\mathcal{P}(\vect{x}):\vect{x}\in K\}.
\]
Since $\widetilde{\mathcal{P}}$ is linear and $\widetilde{\mathcal{P}}([\vect{x}]_{2d})=\mathcal{P}(\vect{x})$, we have
\[
\widetilde{\mathcal{P}}
\big(\mathscr{R}_d(K)\cap\{z_0=1\}\big)
= \operatorname{conv}\big(\mathcal{P}(K)\big).
\]
Thus, passing from $\{[\vect{x}]_{2d}:\vect{x}\in K\}$ to its convex hull corresponds to replacing the generally nonconvex image $\mathcal{P}(K)$ by $\operatorname{conv}(\mathcal{P}(K))$.
}
\rd{For every $\vect{z}\in\mathscr{R}_d(K)\cap\{z_0=1\}$, each quantity $\langle p_i,\vect{z}\rangle$ is nonnegative. 
The logarithmic objective is evaluated only at points for which all these quantities are positive.}
This is because any $\vect{z}\in \mathscr{R}_d(K)\cap \{z_0=1\}$ 
can be decomposed as $\vect{z} = \sum_{j=1}^l \lambda_j[\vect{v}_j]_{2d}$ 
for some $l\in \N$ with all $\vect{v}_j\in K$ and 
$\lambda_j\ge 0,\,\lambda_1+\cdots+\lambda_l=1$.
\rd{By Assumption~\ref{ass:positivity},
$p_i(\vect{v}_j)\geq0$ for every $i\in[m]$ and $j\in[l]$.} Hence,
\[
\langle p_i, \vect{z}\rangle 
= \Big\langle p_i,\sum\limits_{j=1}^l \lambda_j[\vect{v}_j]_{2d}\Big\rangle 
=  \sum\limits_{j=1}^l \lambda_j \langle p_i, [\vect{v}_j]_{2d}\rangle 
= \sum\limits_{j=1}^{l} \lambda_j p_i(\vect{v}_j) \ge 0.
\]
This leads to a natural convex relaxation of \eqref{eq:momequiv}:
\begin{equation}  \label{eq:relax_meas}
	\left\{\begin{array}{rl}
		f_{\operatorname{mom}} \coloneqq \max\limits_{\vect{z}} 
		&  \sum\limits_{i=1}^m a_i \log \langle p_i,\vect{z}\rangle   \\
		\st & z_0 = 1,\, \vect{z}  \in  \mathscr{R}_d(K) .
	\end{array}\right.
\end{equation}
Equivalently, the optimal values of \eqref{eq:problem} and
\eqref{eq:relax_meas} can be written in image space as
\begin{equation}\label{eq:image_formulations}
f_{\max}
=\max_{\vect{u}\in\mathcal{P}(K)}\phi(\vect{u}),
\qquad
f_{\operatorname{mom}}
=\max_{\vect{u}\in\operatorname{conv}(\mathcal{P}(K))}\phi(\vect{u})
\end{equation}
where the maximization is restricted to the domain on which the
logarithms are well defined.
It is clear that $f_{\max}\le f_{\operatorname{mom}}$.
In particular, when $f_{\max}=f_{\operatorname{mom}}$,
the optimization problem \eqref{eq:relax_meas} is said to
be a {\it tight} relaxation of \eqref{eq:problem}.


Assume that both maxima in \eqref{eq:image_formulations} are attained.
Since $\phi$ is strictly concave, the second maximization in
\eqref{eq:image_formulations} has a unique optimizer $\vect{u}^*$. Therefore,
\[
f_{\max}=f_{\operatorname{mom}}
\quad\Longleftrightarrow\quad
\vect{u}^*\in\mathcal{P}(K).
\]
Thus, the convex relaxation \eqref{eq:relax_meas} is tight for
\eqref{eq:problem} if and only if the unique optimizer of
the second maximization in \eqref{eq:image_formulations}
belongs to $\mathcal{P}(K)$. We call
$f_{\operatorname{mom}}-f_{\max}$ the convexification gap.


For a polynomial objective, moment lifting renders the objective
linear in the moment variables, and maximizing a linear function over
a set or over its convex hull gives the same optimal value. In the
present setting, however, the lifted objective is nonlinear and
concave, so maximizing it over the convex hull may yield a strictly
larger optimal value.

{\color{red}
\begin{remark}
The passage from the first image-space problem to the second in
\eqref{eq:image_formulations}, which replaces $\mathcal{P}(K)$ by
$\operatorname{conv}(\mathcal{P}(K))$, is not specific to logarithmic
objectives. More generally, consider an objective of the form $\psi(H(\vect{x}))$, where $H:\mathbb{R}^n\to\mathbb{R}^r$ is a polynomial map and $\psi$ is concave on a convex set containing $\operatorname{conv}(H(K))$.
The corresponding convexification replaces $H(K)$ by $\operatorname{conv}(H(K))$.
Assuming that both optima are attained, tightness holds if and only if the convexified problem has an optimizer belonging to $H(K)$.
When $\psi$ is strictly concave, its optimal point in the image space is unique.

Within this general framework, the logarithmic structure provides two
features specific to the present paper. First, the identity in
\eqref{eq:weighted_product} allows the optimization of a high-degree
polynomial product, when the $a_i$ are positive integers, to be
reformulated using the individual low-degree polynomials.
Second, the resulting relaxation order is independent of the magnitudes of the exponents $a_i$, and the formulation remains available when some exponents are noninteger.

Despite these advantages, replacing the polynomial image by its convex
hull may introduce an arbitrarily large convexification gap
$f_{\operatorname{mom}}-f_{\max}$. This phenomenon is illustrated in
the following example.
\end{remark}
}

\begin{example}\label{ex:gap}
	Consider the univariate log-polynomial optimization problem
	\begin{equation}\label{eq:ex_gap}
		\left\{\begin{array}{cl}
			\max\limits_{x\in\re} & 0.5\log(1+\epsilon-x)+0.5\log(1+\epsilon+x)\\
			\st & 1-x^2 = 0,
		\end{array}
		\right.
	\end{equation}
	where $\epsilon>0$ is a small positive parameter. 
	Since the feasible set $K = \{1,-1\}$ only contains two points, 
	we can express the normalized slice
	$\mathscr{R}_1(K)\cap\{z_0=1\}$ explicitly,
	thus obtain the convex relaxation of \eqref{eq:ex_gap}:
	\[
	\left\{\begin{array}{cl}
		\max\limits_{\vect{z}\in\re^3} & 0.5\log(1+\epsilon-z_1)+0.5\log(1+\epsilon+z_1)\\
		\st &  z_0=z_2 = 1,\, -1\le z_1\le 1.
	\end{array}
	\right.
	\]
	The optimal values of \eqref{eq:ex_gap} and its convex
	relaxation are
	respectively derived as 
	\[
	f_{\max} = 0.5\log(\epsilon)+0.5\log(2+\epsilon),
	\quad
	f_{\operatorname{mom}} = \log(1+\epsilon).
	\]
	The convexification gap satisfies
	\[
	f_{\operatorname{mom}}-f_{\max}
	=
	\log(1+\epsilon)-0.5\log(\epsilon)-0.5\log(2+\epsilon)
	\longrightarrow\infty
	\quad\text{as }\epsilon\to0.
	\]
\end{example}

We next give two concrete sufficient conditions under which the
convexification gap vanishes.

\begin{proposition}\label{prop:gap_maxmom}
Suppose $K$ is nonempty. The convex relaxation
\eqref{eq:relax_meas} is tight for the log-polynomial optimization
\eqref{eq:problem} if either of the following conditions holds.
\begin{enumerate}
\item[(i)] There exists $\vect{\hat{x}}\in K$ that simultaneously maximizes
all the polynomials, i.e.,
\[
p_i(\vect{\hat{x}})=\max_{\vect{x}\in K}p_i(\vect{x})
\qquad\text{for every }i\in[m].
\]

\item[(ii)] The convex relaxation \eqref{eq:relax_meas} has a
maximizer
\[
\vect{z}^*=\sum_{j=1}^l\lambda_j[\vect{v}_j^*]_{2d},
\]
where $\lambda_j\geq0$, $\sum_{j=1}^l\lambda_j=1$, and
$\vect{v}_j^*\in K$, such that
\[
  \mathcal{P}(\vect{v}_1^*)=\cdots=\mathcal{P}(\vect{v}_l^*).
\]
\end{enumerate}
\end{proposition}



\begin{proof} 
The log-polynomial optimization problem \eqref{eq:problem} is equivalent to its moment
reformulation \eqref{eq:momequiv}, and we always have
$f_{\max}\leq f_{\operatorname{mom}}$.

For condition (i), let $\vect{z}$ be any feasible point of
\eqref{eq:relax_meas} and write
\[
\vect{z}=\sum_{j=1}^l\lambda_j[\vect{v}_j]_{2d},
\qquad
\lambda_j\geq0,\quad
\sum_{j=1}^l\lambda_j=1,\quad \vect{v}_j\in K.
\]
For every $i\in[m]$,
\[
\langle p_i,\vect{z}\rangle
=\sum_{j=1}^l\lambda_jp_i(\vect{v}_j)
\leq p_i(\vect{\hat{x}}).
\]
Since the logarithm is increasing,
\[
\sum_{i=1}^m a_i\log\langle p_i,\vect{z}\rangle
\leq
\sum_{i=1}^m a_i\log p_i(\vect{\hat{x}})
\leq f_{\max}.
\]
Taking the maximum over $\vect{z}$ gives
$f_{\operatorname{mom}}\leq f_{\max}$.

For condition (ii), the assumed equality of the polynomial images and
$\sum_{j=1}^l\lambda_j=1$ give
\[
f_{\operatorname{mom}}=\sum_{i=1}^m a_i
  \log\left(\sum_{j=1}^l\lambda_jp_i(\vect{v}_j^*)\right)
=\sum_{i=1}^m a_i\log p_i(\vect{v}_1^*)
\leq f_{\max}.
\]
Thus, under either condition,
$f_{\max}=f_{\operatorname{mom}}$. 
\end{proof}


Condition~(i) requires all the polynomials $p_i$ to attain their
individual maxima at a common point and is therefore a strong special
case. Condition~(ii) has a geometric interpretation. It requires the
optimizer of the second maximization in
\eqref{eq:image_formulations} to belong to $\mathcal{P}(K)$. If both
maximizations in \eqref{eq:image_formulations} attain their optima, then
$f_{\max}=f_{\operatorname{mom}}$ is equivalent to this condition, and
condition~(ii) can be satisfied with $l=1$.


In the special case $m=1$,
the log-polynomial optimization \eqref{eq:problem} is equivalent to 
the polynomial optimization of maximizing $p_1(\vect{x})$ over $K$.
The following result follows directly from Proposition~\ref{prop:gap_maxmom}.
\begin{corollary}
	The relaxation \eqref{eq:relax_meas} is tight if $m=1$.
\end{corollary}

In practice, the problem \eqref{eq:relax_meas} is still difficult to solve.
This is because the moment cone $\mathscr{R}_d(K)$ typically 
does not have a convenient parametric expression in the multivariate case.
On the other hand, the moment cone $\mathscr{R}_d(K)$ 
is often approximated by semidefinite constraints.
This motivates us to further construct 
a hierarchy of semidefinite relaxations of \eqref{eq:relax_meas}.


\subsection{A hierarchy of moment relaxations}

Recall that $K$ is determined by polynomial constraints
\[
c_j(\vect{x}) = 0\,(j\in\mathcal{E}),\quad c_j(\vect{x})\ge 0\,(j\in\mathcal{I}).
\]
Let $\vect{z}\in\mathscr{R}_d(K)$ be an arbitrary truncated moment vector. 
For all $k\ge d$, there exists a $k$-th order truncated multi-sequence extension of $\vect{z}$
such that
\[
\vect{y}=(y_{\vect{\alpha}})\in\re^{\N_{2k}^n}\,\,\mbox{with}\,\,
y_{\vect{\alpha}} = z_{\vect{\alpha}}
\, \forall \vect{\alpha}\in\N_{2d}^n\quad \mbox{and}
\]
\[
\mat{M}_{k}[\vect{y}]\succeq 0,\quad \mat{L}_{c_j}^{(k)}[\vect{y}] = 0\,(j\in \mathcal{E}),\quad
\mat{L}_{c_j}^{(k)}[\vect{y}]\succeq 0\, (j\in \mathcal{I}).
\]
In the above, $\mat{M}_k[\vect{y}]$ is the $k$-th order moment matrix and each
$\mat{L}_{c_j}^{(k)}[\vect{y}]$ is the $k$-th order 
localizing matrix associated with $c_j$, 
which are introduced in Subsection~\ref{ssc:momloc}.
For increasing values of $k = d,d+1,\ldots$, 
we can build a hierarchy of moment relaxations of \eqref{eq:relax_meas} 
with semidefinite constraints: 
\begin{equation}    \label{eq:relax_log}
	\left\{\begin{array}{rl}
		f_{\operatorname{mom},k} \coloneqq \displaystyle \max\limits_{\vect{y}} 
		& \sum\limits_{i=1}^m a_i \log \langle p_i, \vect{y} \rangle \\
		\st  &  y_0 = 1,\,\vect{y} \in \mathbb{R}^{\mathbb{N}_{2k}^n},\, \mat{M}_k [\vect{y}] \succeq 0,\\
		&   \mat{L}_{c_j}^{(k)}[\vect{y}] = 0\,(j \in \mathcal{E}),\\ 
		& \mat{L}_{c_j}^{(k)}[\vect{y}] \succeq 0 \,(j \in \mathcal{I} ),
	\end{array}\right.
\end{equation}
where the integer $k$ is called the {\it relaxation order}.
To distinguish it from the convex relaxation \eqref{eq:relax_meas},
we call the problem \eqref{eq:relax_log} the 
$k$-th order {\it moment relaxation}
and the corresponding hierarchy the {\it moment hierarchy}.
Note that each moment relaxation \eqref{eq:relax_log} is a convex
optimization problem, but it is not a standard semidefinite program
because its objective is concave and nonlinear.
In computational practice, it can be solved globally by conic problem 
solvers with interior point methods.
Since the feasible sets become tighter as the relaxation order
increases, and every point of $\mathscr{R}_d(K)$ admits such extensions,
we have
\[
f_{\operatorname{mom}}\le \cdots\le f_{\operatorname{mom},k+1}
\le f_{\operatorname{mom},k}\le \cdots \le f_{\operatorname{mom},d}.
\]
The following result shows that, under the Archimedean condition, the
moment hierarchy \eqref{eq:relax_log} converges to the convex relaxation
\eqref{eq:relax_meas}. Importantly, this convergence does not by
itself imply convergence to the original optimal value $f_{\max}$.

\begin{theorem}\label{thm:asymp}
Under Assumption~\ref{ass:positivity}, suppose that
\[
\operatorname{Ideal}[\vect{c}_{\mathcal{E}}]
+\operatorname{QM}[\vect{c}_{\mathcal{I}}]
\]
is Archimedean. Then
\[
f_{\operatorname{mom},k}
\longrightarrow f_{\operatorname{mom}}
\qquad\text{as }k\to\infty.
\]
\end{theorem}

{\color{red}
\begin{proof}
Let
\[
C\coloneqq\operatorname{conv}(\mathcal{P}(K)).
\]
The Archimedean condition implies that $K$ is compact. Since
$\mathcal{P}$ is continuous, $\mathcal{P}(K)$ is compact, and so is $C$.

Under the convention $\log 0=-\infty$, the function
\[
\phi(\vect{u})\coloneqq\sum_{i=1}^m a_i\log u_i
\]
is upper semicontinuous on $C$. Moreover, since
$p_i(\vect{x}^{\circ})>0$ for every $i\in[m]$, each
$\log p_i(\vect{x}^{\circ})$ is finite, and hence
$\phi(\mathcal{P}(\vect{x}^{\circ}))
=\sum_{i=1}^m a_i\log p_i(\vect{x}^{\circ})$ is finite.
It follows that $\phi$ attains its maximum over $C$ at
some point $\vect{u}^*$, and necessarily
\[
u_i^*>0\qquad(i\in[m]).
\]

Since $\vect{u}^*$ maximizes the differentiable concave function $\phi$
over the convex set $C$, the first-order optimality condition gives
\[
\nabla\phi(\vect{u}^*)^T(\vect{u}-\vect{u}^*)\leq0
\qquad\text{for every }\vect{u}\in C.
\]
In particular, the polynomial
\[
q(\vect{x})\coloneqq
-\nabla\phi(\vect{u}^*)^T\big(\mathcal{P}(\vect{x})-\vect{u}^*\big)
=
\sum_{i=1}^m\frac{a_i}{u_i^*}
\big(u_i^*-p_i(\vect{x})\big)
\]
is nonnegative on $K$.

For every $\epsilon>0$, the polynomial $q+\epsilon$ is strictly
positive on $K$. By Putinar's Positivstellensatz \cite{Putinar}, there
exists an
integer $N(\epsilon)\geq d$ such that
\[
q+\epsilon\in
\operatorname{Ideal}[\vect{c}_{\mathcal{E}}]_{2k}
+\operatorname{QM}[\vect{c}_{\mathcal{I}}]_{2k}
\qquad
\text{for every }k\geq N(\epsilon).
\]

Fix $k\geq N(\epsilon)$ and let
$\vect{y}\in\mathbb{R}^{\mathbb{N}_{2k}^n}$ be a feasible point of
\eqref{eq:relax_log} satisfying
$\langle p_i,\vect{y}\rangle>0$ for every $i\in[m]$. Define
\[
\vect{u}(\vect{y})\coloneqq
\big(\langle p_1,\vect{y}\rangle,\ldots,\langle p_m,\vect{y}\rangle\big).
\]
Since $\deg(p_i)\leq2d\leq2k$ for every $i\in[m]$, the vector $\vect{u}(\vect{y})$
depends only on the degree-$2d$ truncation of $\vect{y}$.
The certificate above and the feasibility of $\vect{y}$ give
\[
\langle q+\epsilon,\vect{y}\rangle\geq0.
\]
Indeed, the ideal terms pair to zero, while the quadratic-module
terms pair nonnegatively because the corresponding localizing matrices
of $\vect{y}$ are positive semidefinite.
By the definition and linearity of $\langle\cdot,\vect{y}\rangle$ in
\eqref{eq:bilinear}, and since $\langle1,\vect{y}\rangle=y_0=1$, we have
\begin{align*}
\langle q+\epsilon,\vect{y}\rangle
&=
\sum_{i=1}^m\frac{a_i}{u_i^*}
\left(
u_i^*\langle1,\vect{y}\rangle-\langle p_i,\vect{y}\rangle
\right)
+\epsilon\langle1,\vect{y}\rangle\\
&=
\sum_{i=1}^m\frac{a_i}{u_i^*}
\left(
u_i^*-\langle p_i,\vect{y}\rangle
\right)
+\epsilon\\
&=
-\nabla\phi(\vect{u}^*)^T\big(\vect{u}(\vect{y})-\vect{u}^*\big)+\epsilon.
\end{align*}
Combining this identity with
$\langle q+\epsilon,\vect{y}\rangle\geq0$ yields
\[
\nabla\phi(\vect{u}^*)^T\big(\vect{u}(\vect{y})-\vect{u}^*\big)\leq\epsilon.
\]
By concavity of $\phi$,
\[
\phi(\vect{u}(\vect{y}))
\leq
\phi(\vect{u}^*)+
\nabla\phi(\vect{u}^*)^T\big(\vect{u}(\vect{y})-\vect{u}^*\big)
\leq
\phi(\vect{u}^*)+\epsilon
=
f_{\operatorname{mom}}+\epsilon.
\]
Taking the supremum over all feasible $\vect{y}$ yields
\[
f_{\operatorname{mom},k}
\leq f_{\operatorname{mom}}+\epsilon
\qquad
\text{for every }k\geq N(\epsilon).
\]
Since $f_{\operatorname{mom}}\leq f_{\operatorname{mom},k}$, we have
\[
f_{\operatorname{mom}}
\leq f_{\operatorname{mom},k}
\leq f_{\operatorname{mom}}+\epsilon
\qquad
\text{for every }k\geq N(\epsilon).
\]
Therefore,
\[
f_{\operatorname{mom},k}
\longrightarrow f_{\operatorname{mom}}
\qquad\text{as }k\to\infty.
\]
\end{proof}
}

Theorem~\ref{thm:asymp} concerns only the moment-truncation gap.
Convergence to the original optimal value additionally requires the
convexification gap to vanish.

\begin{corollary}\label{cor:asymp}
Assume that
$\operatorname{Ideal}[\vect{c}_{\mathcal{E}}]
+\operatorname{QM}[\vect{c}_{\mathcal{I}}]$
is Archimedean. If
$f_{\max}=f_{\operatorname{mom}},$
then
$f_{\operatorname{mom},k}\longrightarrow f_{\max}\text{ as }k\to\infty.$
Since either condition in Proposition~\ref{prop:gap_maxmom} implies
$f_{\max}=f_{\operatorname{mom}}$, the same convergence conclusion
holds whenever either condition is satisfied.
\end{corollary}


\begin{proof}
The conclusion follows immediately from
Theorem~\ref{thm:asymp} and
Proposition~\ref{prop:gap_maxmom}.
\end{proof}




\section{Tightness analysis of moment relaxations}\label{sec:tight}

\rd{The $k$-th moment relaxation \eqref{eq:relax_log} is called
{\it tight} if its optimal value equals $f_{\max}$. The moment hierarchy
is said to have {\it finite convergence} if this equality holds at some
finite relaxation order.}
In this section, we study how to certify such tightness and 
\rd{how to extract global maximizers when an appropriate flat-truncation
condition holds.}
\rd{We also study a class of problems with SOS-concave polynomials in
the logarithmic objective, SOS-concave inequality polynomials, and
linear equality constraints, for which finite convergence is guaranteed.}


\subsection{Certifying tightness of moment relaxations}

Let $k$ be the relaxation order.
Let $f_{\operatorname{mom},k}$ denote the optimal value
of the moment relaxation \eqref{eq:relax_log},
$f_{\operatorname{mom}}$ the optimal value
of the convex relaxation \eqref{eq:relax_meas},
and $f_{\max}$ the optimal value of the
original log-polynomial optimization \eqref{eq:problem}.
Since
\[
f_{\max}\leq f_{\operatorname{mom}}
\leq f_{\operatorname{mom},k},
\]
\rd{the $k$-th moment relaxation is tight if and only if the following
two equalities hold:}
\[
\underbrace{f_{\max}=f_{\operatorname{mom}}}_{\text{no convexification gap}},
\qquad
\underbrace{f_{\operatorname{mom}}
=f_{\operatorname{mom},k}}_{\text{no moment-truncation gap}}
\]
\rd{The first equality asks whether replacing $\mathcal{P}(K)$ by
$\operatorname{conv}(\mathcal{P}(K))$ in
\eqref{eq:image_formulations} changes the optimal value. Under
attainment, this is equivalent to asking whether the second maximization
in \eqref{eq:image_formulations} has an optimizer belonging to
$\mathcal{P}(K)$, as discussed in Section~\ref{sec:rel}. Let
$\vect{y}^*$ be an optimizer of the $k$-th moment relaxation
\eqref{eq:relax_log}. The flat truncation condition below certifies the
second equality by showing that $\vect{y}^*|_{2t}$ admits a
$K$-representing measure.}
It reads
\begin{equation}\label{eq:flat_truncation}
	\exists t\in\rd{\{d,d+1,\ldots,k\}}\quad \mbox{s.t.}\quad
	\rank\,\mat{M}_t[\vect{y}^*] = \rank\,\mat{M}_{t-d}[\vect{y}^*].
\end{equation}
Under the flat truncation condition, \rd{the truncation
$\vect{y}^*|_{2t}$} admits a unique $K$-representing
atomic measure supported by $r=\rank\, \mat{M}_t[\vect{y}^*]$ distinct points.
In other words, \rd{there exist scalars
$\lmd_1,\ldots,\lmd_r>0$}
and distinct points $\vect{v}_1^* \ddd \vect{v}_r^* \in K$ such that
\begin{equation}\label{eq:ft_decomp}
	\vect{y}^*|_{2t} = \lmd_1 [\vect{v}_1^*]_{2t} + \cdots + \lmd_r [\vect{v}_r^*]_{2t},
\end{equation}
where $\vect{y}^*|_{2t}$ is \rd{the degree-$2t$ truncation} of $\vect{y}^*$.
\rd{Since $y_0^*=1$, the weights satisfy
$\sum_{j=1}^r\lambda_j=1$.}
In particular, such a decomposition \eqref{eq:ft_decomp} is unique by 
\cite[Theorem~2.7.7]{MPO}, up to reordering. 
Since $t\ge d$, the decomposition \eqref{eq:ft_decomp} ensures that 
$\vect{y}^*|_{2d}\in \mathscr{R}_d(K)$.
We then summarize sufficient conditions for \eqref{eq:relax_log} 
to be a tight relaxation for \eqref{eq:problem}.


\begin{theorem}\label{tm:parallel_all}
	Suppose $\vect{y}^*$ is an optimizer of the $k$-th order moment relaxation \eqref{eq:relax_log}.
	If $\vect{y}^*$ satisfies the flat truncation condition 
	\eqref{eq:flat_truncation} for some
	$t\in\rd{\{d,d+1,\ldots,k\}}$,
	then $f_{\operatorname{mom},k} = f_{\operatorname{mom}}$,
	and $\vect{y}^*|_{2t}$ admits the decomposition in \eqref{eq:ft_decomp} 
	for $r=\rank\, \mat{M}_t[\vect{y}^*]$ distinct points $\vect{v}_1^*,\ldots, \vect{v}_r^*\in K$.
	If, in addition,
$\mathcal{P}(\vect{v}_1^*)=\cdots=\mathcal{P}(\vect{v}_r^*)$,
	then $f_{\max} = f_{\operatorname{mom}} = f_{\operatorname{mom},k}$
	and all $\vect{v}_1^* \ddd \vect{v}_r^*$ are global maximizers of \reff{eq:problem}.
\end{theorem}

\begin{proof} 
	Under the flat truncation condition, 
	$\vect{z}^* = \vect{y}^*|_{2d}$ is a feasible point of \eqref{eq:relax_meas}.
	Since each $\deg(p_i) \le 2d$, we have
	\begin{equation}
		f_{\operatorname{mom},k} = \sum_{i=1}^m a_i \log \langle p_i,\vect{y}^*\rangle 
		= \sum_{i=1}^m a_i \log \langle p_i,\vect{z}^*\rangle 
		\leq f_{\operatorname{mom}}.
	\end{equation}
	\rd{Combining} the above inequality with
	$f_{\operatorname{mom}} \le f_{\operatorname{mom},k}$,
	we conclude that $f_{\operatorname{mom}} = f_{\operatorname{mom},k}$.
	The decomposition \eqref{eq:ft_decomp} is implied by the 
	flat truncation condition \cite[Theorem~2.7.7]{MPO}.
	If
$\mathcal{P}(\vect{v}_1^*)=\cdots=\mathcal{P}(\vect{v}_r^*)$,
	then $f_{\max} = f_{\operatorname{mom}}$ by Proposition~\ref{prop:gap_maxmom}.
	Consequently, $f_{\max} = f_{\operatorname{mom}} = f_{\operatorname{mom},k}$.
	To show that $\vect{v}_l^*$ for each $l\in[r]$ is a global maximizer, 
	it suffices to show $f(\vect{v}_l^*) = f_{\operatorname{mom},k} = f_{\max}$.
	\rd{Using
	$\mathcal{P}(\vect{v}_1^*)=\cdots=\mathcal{P}(\vect{v}_r^*)$
	and $\sum_{j=1}^r\lambda_j=1$,} we have
	\begin{equation*}
		f_{\operatorname{mom},k}
		= \sum_{i=1}^m a_i \log \Big( \sum_{j=1}^r \lambda_j p_i(\vect{v}_j^*) \Big) 
		= \sum_{i=1}^m a_i \log \big( p_i(\vect{v}_l^*) \big) = f(\vect{v}_l^*)
	\end{equation*}
	for every $l\in[r]$. It follows that $\vect{v}_1^*, \ldots, \vect{v}_r^*$ 
	are all global maximizers of \eqref{eq:problem}. 
\end{proof}

Notice that $\rank\, \mat{M}_d[\vect{y}^*] = 1$ if and only if $\vect{y}^*|_{2d} = [\vect{v}^*]_{2d}$ 
with $\vect{v}^* = (y^*_{\vect{e}_1}\ddd y^*_{\vect{e}_n})$. 
For this special case, all conditions in Theorem~\ref{tm:parallel_all} 
are satisfied automatically.

\begin{corollary}\label{cor:rank1}
	Suppose $\vect{y}^*$ is an optimizer of the $k$-th order moment relaxation
	\eqref{eq:relax_log}.
	If $\rank\, \mat{M}_d[\vect{y}^*] = 1$, then $f_{\max} = f_{\operatorname{mom},k}$
	and $\vect{v}^* = (y^*_{\vect{e}_1}\ddd y^*_{\vect{e}_n})$
	is a global maximizer of \eqref{eq:problem}.
\end{corollary}

\rd{The role of flat truncation here is closely related to its role
in standard polynomial optimization \cite{nie2013certifying}.}
Suppose $\vect{y}^*$ is an optimizer of the $k$-th order moment relaxation
\eqref{eq:relax_log}.
If it satisfies the flat truncation condition
\eqref{eq:flat_truncation}, then \eqref{eq:relax_log} is a tight
relaxation of \eqref{eq:relax_meas}, which implies that
\[
f_{\operatorname{mom}}=f_{\operatorname{mom},k}.
\]
\rd{Thus, flat truncation certifies that the moment-truncation gap vanishes but does
not by itself certify that the convexification gap vanishes.}
The convex relaxation \eqref{eq:relax_meas} may not be a tight
relaxation of the moment reformulation \eqref{eq:momequiv}, as shown
in Example~\ref{ex:gap}.
Consequently, even if $\vect{y}^*|_{2t}$ admits the decomposition
\eqref{eq:ft_decomp} with candidate solutions
$\vect{v}_1^*,\ldots,\vect{v}_r^*\in K$, it remains possible that
\[
f_{\max}<f_{\operatorname{mom}}
=f_{\operatorname{mom},k}.
\]
In such scenarios, the points $\vect{v}_i^*$ may or may not be global
maximizers of \eqref{eq:problem}.
We illustrate both possibilities in the following example.



\begin{example}\label{ex:counter}
	Consider the log-polynomial optimization problem:
	\begin{equation}\label{eq:counter}
		\left\{\begin{array}{cl}
			\max\limits_{x\in\re} & f(x)\coloneqq \log((x-1)^2+a) + \log((x+1)^2+a)\\
			\st  & x(x-1)(x+1) = 0,
		\end{array}\right.
	\end{equation}
	where $a>0$ is a parameter. The feasible set $K = \{1, -1, 0\}$
	and the degree bound $d=2$.
	It is easy to evaluate
	\[ f(1) = \log (a(4+a)), \quad f(-1) = \log (a(4+a)),\quad  f(0) = 2\log (1+a). \]
	Consider the moment relaxation \eqref{eq:relax_log} with 
	the relaxation order $k=3$.
	\rd{The equality constraint in \eqref{eq:counter} is equivalent to
	$x^3-x=0$. Under the localizing-matrix definition used in
	\eqref{eq:relax_log}, it yields
	$y_3=y_1$, $y_4=y_2$, and $y_5=y_3$, while $y_6$ remains a
	free moment variable.}
	By eliminating variables, 
	the $3$rd-order moment relaxation of \eqref{eq:counter} can be simplified to 
	\begin{equation}\label{eq:relax_counter}
		\left\{\begin{array}{cl}
			\displaystyle \max_{\rd{(y_1,y_2,y_6)}} &
			\log(y_2-2y_1+1+a ) +  \log(y_2+2y_1+1+a )   \\
			\st &
			\rd{y_1,y_2,y_6\in\re},\,\left[ \begin{array}{cccc}
				1 & y_1 & y_2 & y_1 \\
				y_1 & y_2 & y_1 & y_2\\
				y_2 & y_1 & y_2 & y_1\\
				y_1 & y_2 & y_1 & \rd{y_6}\\
			\end{array}
			\right] \succeq 0.
		\end{array}\right.
	\end{equation}
	In the above, the objective function equals $\log((y_2+1+a)^2-4y_1^2)$ 
	and the positive semidefinite constraint implies that 
	$0\le y_1^2 \le y_2^2 \le y_2 \le 1.$
	Thus, \rd{for every feasible point of \eqref{eq:relax_counter}, we have}
	\[
	\log((y_2+1+a)^2-4y_1^2)\le \log((1+1+a)^2-4(0)) = 2\log(a+2).
	\]
	\rd{For $y_1=0$ and $y_2=1$, the positive semidefinite constraint in
	\eqref{eq:relax_counter} holds if and only if $y_6\geq1$.
	Since the objective is independent of $y_6$, every such choice attains
	the above upper bound. Choosing $y_6=1$ gives a flat optimal solution.}
	Hence the optimal value and \rd{one flat optimizer} of the 
	$3$rd-order moment relaxation of \eqref{eq:counter} are
	\[
	f_{\operatorname{mom},3} = 2\log(a+2),\quad
	\vect{y}^* = (1,0,1,0,1,0,1),
	\] 
	where $\vect{y}^* = 0.5[u_1]_6 + 0.5[u_2]_6$ is supported 
	on feasible points $u_1 = 1$, $u_2 = -1$.
	
	(i) When $a = 1$, we have
	\[
	f(1) = f(-1) = \log(5)>f(0) = \log(4).
	\]
	In this case, $f_{\max} = \log(5)$ and 
	both $u_1 = 1$ and $u_2 = -1$ are maximizers of (\ref{eq:counter}).
	However, $f_{\operatorname{mom},3} = \log(9)$, 
	which is strictly bigger than $f_{\max}$.
	Indeed, one can check that 
	$f_{\operatorname{mom}} = f_{\operatorname{mom},k} = \log(9)$ 
	for all $k\ge 3$, \rd{indicating a persistent nonzero convexification gap.}
	
	(ii) When $a = 0.1$, we have
	\[
	f(1) =  f(-1) = \log (0.41)< f(0) = \log (1.21).
	\]
	This implies that $f_{\max} = \log(1.21)$,
	which is achieved at the unique optimizer $u_3 = 0$.
	In this case, $f_{\operatorname{mom},3} = \log (4.41)>f_{\max}$ 
	and $u_1,u_2$ are no longer global maximizers of \eqref{eq:counter}.
	It can be further verified that the relaxation \eqref{eq:relax_log} 
	remains non-tight for all $k\ge 3$.
\end{example}

Proposition~\ref{prop:gap_maxmom} gives two sufficient conditions for 
$f_{\max} = f_{\operatorname{mom}}$.
These conditions are problem-dependent and may be difficult to
verify a priori.
Recall that moment relaxations for polynomial optimization exhibit
finite convergence under suitable optimality conditions
\cite{nie2014optimality}. This motivates us to study a convex regime
of log-polynomial optimization.



\subsection{Concave log-polynomial optimization}

The problem \eqref{eq:problem} is called a {\it concave} log-polynomial 
optimization problem if $K$ is a convex set and 
each $p_i$ is concave on $K$, i.e., 
for all $\vect{u},\vect{v}\in K$, $i\in[m]$ and $\lambda\in[0,1]$,
\[
\lambda \vect{u}+(1-\lambda)\vect{v}\in K,\quad
p_i(\lambda \vect{u}+(1-\lambda)\vect{v})\ge  \lambda p_i(\vect{u})+(1-\lambda)p_i(\vect{v}).
\]
Concave log-polynomial optimization has tight moment relaxations 
under flat truncation conditions.

\begin{theorem}   \label{thm:concave}
	Assume \eqref{eq:problem} is concave and 
	$\vect{y}^*$ is an optimizer of its $k$-th order moment relaxation \reff{eq:relax_log}.
	If $\vect{y}^*$ satisfies the flat truncation condition \reff{eq:flat_truncation},
	then $f_{\max} = f_{\operatorname{mom}} = f_{\operatorname{mom},k}$ 
	and $\vect{v}^* = (y^*_{\vect{e}_1}\ddd y^*_{\vect{e}_n})$
	is a global maximizer of \reff{eq:problem}.
\end{theorem}

\begin{proof} 
	Suppose $\vect{y}^*$ satisfies the flat truncation condition \eqref{eq:flat_truncation}
	for some $t\in\{d,d+1,\ldots,k\}$. Let
	$r = \rank\,\mat{M}_t[\vect{y}^*]$.
	There exist positive scalars $\lambda_i$ and distinct points $\vect{v}_i^*\in K$
	for $i\in[r]$ such that 
	\[
	\vect{y}^*|_{2t} = \lambda_1[\vect{v}_1^*]_{2t}+\cdots
	+\lambda_r[\vect{v}_r^*]_{2t}.
	\]
	\rd{The equality $y_0^*=1$ implies
	$\lambda_1+\cdots+\lambda_r=1$. Since $K$ is convex,}
	$\vect{v}^* = (y_{\vect{e}_1}^*, \ldots, y_{\vect{e}_n}^*) 
	= \lambda_1\vect{v}_1^*+\cdots+\lambda_r\vect{v}_r^*\in K.$
	\rd{Since $\deg(p_i)\leq2d\leq2t$, this implies}
	\[
	\langle p_i,\vect{y}^*\rangle 
	= \sum_{j=1}^r \lambda_j \langle p_i, \rd{[\vect{v}_j^*]_{2t}}\rangle 
	= \sum_{j=1}^r \lambda_j p_i(\vect{v}_j^*) \le p_i(\vect{v}^*)
	\]
	for all $i\in[m]$.
	By assumption, $\vect{y}^*$ is an optimizer of \eqref{eq:relax_log}.
	Since the logarithm function is monotonically increasing, 
	we further have
	\[
	f_{\operatorname{mom},k}
	= \sum\limits_{i=1}^m a_i \log \langle p_i, \vect{y}^* \rangle  
	\le \sum\limits_{i=1}^m a_i \log p_i(\vect{v}^*)
	\le f_{\max}.
	\]
	\rd{Since \eqref{eq:relax_meas} and \eqref{eq:relax_log} are
	relaxations of \eqref{eq:problem}, we also have
	$f_{\max}\leq f_{\operatorname{mom}}\leq
	f_{\operatorname{mom},k}$. Consequently,
	$f_{\max}=f_{\operatorname{mom}}=f_{\operatorname{mom},k}$, and
	$\vect{v}^*$ is a maximizer of \eqref{eq:problem}.}
\end{proof}

\rd{Theorems~\ref{tm:parallel_all} and \ref{thm:concave} provide two
different sufficient mechanisms for obtaining tight moment
relaxations, both under flat truncation.}
Specifically, Theorem~\ref{tm:parallel_all} provides conditions 
under which all points $\vect{v}_1^*, \ldots, \vect{v}_r^*$ 
in the support of the atomic measure admitted by $\vect{y}^*|_{2t}$ 
are maximizers of \eqref{eq:problem}.
In contrast, Theorem~\ref{thm:concave} focuses on concave log-polynomial optimization, 
where these $\vect{v}_1^*, \ldots, \vect{v}_r^*$ are not necessarily 
global maximizers of \eqref{eq:problem}.

However, determining the convexity or concavity of polynomials 
with degree four or higher is NP-hard \cite{AAA13}. 
In computational practice, SOS-convex/concave polynomials, 
which form a subset of convex/concave polynomials, 
are of great interest since they can be verified by
solving a semidefinite program \cite[Lemma~7.1.3]{MPO}.
A polynomial $q(\vect{x})$ is said to be {\it SOS-convex} 
if there exists a matrix polynomial $\mat{R}(\vect{x})$ such that its Hessian matrix
$\nabla^2 q(\vect{x})
=\mat{R}(\vect{x})^T\mat{R}(\vect{x})$,
and it is said to be {\it SOS-concave} if $-q$ is SOS-convex.
We then show that $f_{\max} = f_{\operatorname{mom},k}$ for every $k\ge d$
if \eqref{eq:problem} is defined by SOS-concave polynomials.


\begin{theorem}\label{thm:sosconcave}
	Assume $c_j$ are SOS-concave for all $j\in\mathcal{I}$, 
	$\vect{c}_{\mathcal{E}}$ is a linear polynomial tuple and 
	each $p_i$ for $i\in[m]$ is SOS-concave.
	Then $f_{\operatorname{mom},k} = f_{\max}$ for all $k\ge d$.
	\rd{Moreover, for any $k\geq d$, if the $k$-th order moment
	relaxation \eqref{eq:relax_log} has an optimizer $\vect{y}^*$,}
	then $\vect{v}^* = (y^*_{\vect{e}_1}\ddd y^*_{\vect{e}_n})$
	is a global maximizer of \eqref{eq:problem}.
\end{theorem}

\begin{proof} 
	Given $k\ge d$, let $\vect{y}$ be a feasible point of \eqref{eq:relax_log}
	and denote $\pi(\vect{y})\coloneqq (y_{\vect{e}_1}\ddd y_{\vect{e}_n})$.
	For every $j\in \mathcal{I}$, $\langle c_j,\vect{y}\rangle \ge 0$ 
	since $\mat{L}_{c_j}^{(k)}[\vect{y}]\succeq 0$.
	For every $j\in \mathcal{E}$, since $c_j$ is linear, 
	we have $c_j(\pi(\vect{y})) = 0$ as it corresponds to the 
	$(1,1)$-th entry of $\mat{L}_{c_j}^{(k)}[\vect{y}]$.
	Since $\mat{M}_d[\vect{y}]\succeq 0$ and $y_0=1$, 
	the SOS-concavity and Jensen's inequality \cite[Theorem~7.1.6]{MPO} imply
	\[ 
	\langle p_i,\vect{y}\rangle \le p_i(\pi(\vect{y})),\quad 0
	\le \langle c_j,\vect{y}\rangle \le c_j(\pi(\vect{y})), 
	\]
	for all $i\in[m]$ and $j\in \mathcal{I}$. 
	Thus, we have $\pi(\vect{y})\in K$ and
	\begin{equation}\label{eq:jensen}
		\sum\limits_{i=1}^m a_i \log \langle p_i,\vect{y}\rangle  
		\le  \sum\limits_{i=1}^m a_i \log p_i(\pi(\vect{y})) 
		\leq  f_{\max}.
	\end{equation}
	Note that the relaxation order $k$ and the 
	feasible point $\vect{y}$ can be chosen arbitrarily.
	By maximizing \eqref{eq:jensen} with respect to $\vect{y}$ 
	over the feasible set of \eqref{eq:relax_log}, 
	we obtain $f_{\operatorname{mom},k} \le f_{\max}$, 
	The reverse inequality
	$f_{\max}\leq f_{\operatorname{mom},k}$ follows because
	\eqref{eq:relax_log} is a relaxation of \eqref{eq:problem}.
	Hence $f_{\operatorname{mom},k}=f_{\max}$ for all $k\geq d$.
	Suppose \reff{eq:relax_log} has an optimizer $\vect{y}^*$. 
    Since $\vect{y}^*$ is optimal, the left-hand side of \eqref{eq:jensen} equals
    $f_{\operatorname{mom},k}=f_{\max}$.
    Hence $\sum_{i=1}^m a_i \log p_i(\pi(\vect{y}^*))=f_{\max}$, and $\vect{v}^*=\pi(\vect{y}^*)$ is a global maximizer of \reff{eq:problem}. 
\end{proof}

\rd{Since the logarithm is increasing and concave, if the polynomials
$p_i$ are concave and $K$ is convex, then the log-polynomial
objective is concave on its effective domain. Hence, maximizing it
over $K$ is a convex optimization problem and is amenable to
standard convex optimization methods. The main significance of
Theorem~\ref{thm:sosconcave} in this setting is therefore the
finite-order exactness of the moment relaxation and, when a moment
optimizer exists, the recovery of a global maximizer from its first
moments.}

\rd{Theorem~\ref{thm:sosconcave} provides a finite-order moment
certificate for maximizing a product of positive SOS-concave
polynomials over a convex feasible set.}
Consider a polynomial optimization problem of the form
\begin{equation}\label{eq:prod_pop}
	\left\{\begin{array}{cll}
		\max\limits_{\vect{x}\in\re^n} & \theta(\vect{x})\coloneqq p_1(\vect{x})p_2(\vect{x})\cdots p_{m}(\vect{x})\\
		\st & c_j(\vect{x}) = 0 \,\, (j \in \mathcal{E}),\\
		& c_j(\vect{x}) \geq 0 \,\, (j \in \mathcal{I}).
	\end{array}\right.
\end{equation}
Let $K$ denote the feasible set of \eqref{eq:prod_pop}.
Assume $p_1\ddd p_m$ (not necessarily distinct) are SOS-concave,
each $c_{j}$ for $j\in\mathcal{E}$ is linear, and every $c_{j}$ for $j\in\mathcal{I}$ is SOS-concave.
When $m > 1$, its objective function $\theta(\vect{x})$ is typically not concave.
\rd{Consequently, convexity-based exactness guarantees for the standard
Moment--SOS hierarchy do not generally apply, and its lowest-order
relaxation need not be tight.}
If each $p_i(\vect{x})$ for $i\in[m]$ is positive on $K$, then \eqref{eq:prod_pop} can 
be equivalently solved as the log-polynomial optimization problem
\[
\rd{\max\limits_{\vect{x}\in K}\log \theta(\vect{x})
=\max\limits_{\vect{x}\in K}\sum\limits_{i=1}^m\log p_i(\vect{x}).}
\]
By Theorem~\ref{thm:sosconcave}, this reformulated problem has 
a tight moment relaxation at its lowest relaxation order.
\rd{Beyond this convex regime, genuinely nonconcave log-polynomial
objectives are examined numerically in Section~\ref{sec:scaling}.}




\section{Tighter relaxations with LMEs}
\label{sec:LME}


In this section, we introduce a technique called Lagrange multiplier
expressions to construct tighter moment relaxations of
\eqref{eq:problem} than the standard moment relaxation
\eqref{eq:relax_log}.

For convenience, suppose that $\mathcal{E}\cup \mathcal{I} = [\ell] \coloneq \{1,\ldots, \ell\}$.
Let $f(\vect{x})$ be the objective function and $\vect{x}^*$ the maximizer of \eqref{eq:problem}.
Under a suitable constraint qualification at $\vect{x}^*$, there
exists a vector of Lagrange multipliers
$\vect{\lambda} = (\lambda_j)$ for $j = 1\ddd \ell$, 
such that $(\vect{x}^*,\vect{\lambda})\in\mathcal{K}$, where
\begin{equation}    \label{eq:KKT}
	\mathcal{K} \coloneqq 
	\left\{ (\vect{x}, \vect{\lambda})\in K\times \re^{\ell}
	\left|\begin{array}{c}
		\nabla f(\vect{x}) + \sum\limits_{j \in \mathcal{E} \cup \mathcal{I}} \lambda_j \nabla c_j(\vect{x}) = 0,\\
		\lambda_j\ge 0,\, \lambda_j c_j(\vect{x}) = 0\, (j \in \mathcal{I})
	\end{array}\right.\right\}.
\end{equation}
The constraints in \eqref{eq:KKT} are called the 
Karush-Kuhn-Tucker (KKT) optimality conditions.
Each $(\vect{x}, \vect{\lambda}) \in \mathcal{K}$ is called a critical pair, 
where $\vect{x}$ alone is called a critical point.
Suppose there exists a polynomial or rational tuple 
$\vect{\tau} = (\tau_1,\ldots, \tau_{\ell})$ such that 
\begin{equation}\label{eq:lme}
	\lambda_j = \tau_j(\vect{x})\quad \mbox{for all}\quad
	(\vect{x},\vect{\lambda})\in \mathcal{K}.
\end{equation}
Then we can get a parametric approximation of critical points by
replacing $\vect{\lambda}$ in \eqref{eq:KKT} with $\vect{\tau}$.
Such a $\vect{\tau}$ is called a Lagrange multiplier expression (LME).
LMEs serve as an efficient tool for constructing tight relaxations 
in polynomial optimization \cite{Nie2019}.
Recently, they have been applied to challenging optimization problems 
given by polynomials, including bilevel optimization, 
generalized Nash equilibrium problems and variational inequalities
\cite{Choi2024,NieSunTangZ,nie2023convex,NieTangZgnep21,NieWangZBilevel21}.

\begin{example}\label{ex:lmes}
	We present explicit LMEs for box, simplex and ball constraints.
	
	\begin{enumerate}
		
		\item[(i)] If $K=\{\vect{x}\in\re^n: a_j \leq x_j \leq b_j,\, j\in[n]\}$, 
		then $\ell = 2n$ and \eqref{eq:lme} is satisfied for 
		$\vect{\tau} = (\tau_1,\ldots,\tau_{2n})$ with
		\begin{equation*}
			\tau_{2j-1} = \rd{\frac{\frac{\partial f}{\partial x_j}(\vect{x})(x_j-b_j)}
			{b_j-a_j}},\qquad
			\tau_{2j} = \rd{\frac{\frac{\partial f}{\partial x_j}(\vect{x})(x_j-a_j)}
			{b_j-a_j}}\quad
			\forall j\in[n].
		\end{equation*}
		
		\item[(ii)] If $K=\{\vect{x}\in\re^n: 1 - \mathbf{1}^T \vect{x} \ge 0,\vect{x} \geq 0\}$,
		then $\ell = n+1$ and \eqref{eq:lme} is satisfied for 
		\[
		\vect{\tau} = \vect{x}^T\nabla f(\vect{x})\mathbf{1}
		- (0,\, \nabla_{x_1} f(\vect{x}),
		\, \ldots,\, \nabla_{x_n} f(\vect{x})),
		\]
		where $\mathbf{1}\in\re^{\rd{n+1}}$ is the vector of all ones.
		
		\item[(iii)] If $K = \{\vect{x}\in\re^n: \|\vect{x}-\vect{a}\| \leq b\}$, 
		then $\ell=1$ and \eqref{eq:lme} is satisfied for
		\begin{equation*}
			\vect{\tau}
			= \frac{(\vect{x}-\vect{a})^T \nabla f(\vect{x})}{2b^2}.
		\end{equation*}
		
	\end{enumerate}
\end{example}

Next, we introduce how to apply LMEs to construct tighter
moment relaxations for \eqref{eq:problem}.
Consider the log-polynomial optimization problem (\ref{eq:problem})
with box, simplex or ball constraints.
Then, polynomial LMEs $\tau_j(\vect{x})$ for $j\in\mathcal{E}\cup \mathcal{I}$ are 
explicitly presented in Example~\ref{ex:lmes}. Recall that 
\[
f(\vect{x}) = \sum\limits_{i=1}^m a_i \log p_i(\vect{x}),\quad
\nabla f(\vect{x}) = \sum\limits_{i=1}^m \frac{a_i}{p_i(\vect{x})}\nabla p_i(\vect{x}).
\]
For convenience, denote $p^{\mathbf{1}} = p_1p_2\cdots p_m$ and
\[
f_i(\vect{x}) \coloneqq \frac{p^{\mathbf{1}}(\vect{x})}{p_i(\vect{x})},\quad \forall i\in[m].
\]
Assume \rd{$p_i>0$ on $K$ for all $i\in[m]$}.
Then a feasible point $\vect{x}\in K$ is a critical point of \eqref{eq:problem} 
if and only if it satisfies 
\[
\left\{\begin{array}{l}
	\sum\limits_{i=1}^m a_i f_i(\vect{x})\nabla p_i(\vect{x}) 
	+ p^{\mathbf{1}}(\vect{x})\sum\limits_{j\in \mathcal{E}\cup \mathcal{I}} 
	\tau_j(\vect{x})\nabla c_j(\vect{x}) = 0,\\
	\tau_j(\vect{x})\ge 0\,(j\in\mathcal{I}),\quad  \tau_j(\vect{x})c_j(\vect{x})=0\, (j\in\mathcal{I}).
\end{array}\right.
\]
This system provides additional polynomial constraints for
strengthening the moment relaxation of \eqref{eq:problem}.
For convenience, denote the polynomial tuples
\[	\begin{array}{l}
	\Phi\coloneq \{ \vect{c}_{\mathcal{E}} \} \cup \Big\{ \sum\limits_{i=1}^m 
	a_i f_i(\vect{x})\nabla p_i(\vect{x}) + p^{\mathbf{1}}(\vect{x})\sum\limits_{j\in \mathcal{E}\cup \mathcal{I}} 
	\tau_j(\vect{x})\nabla c_j(\vect{x}) \Big\} \cup 
	\{\tau_jc_j: j\in \mathcal{I}\},  \\
	\Psi\coloneq \{ \vect{c}_{\mathcal{I}} \} \cup \{ \tau_j: j\in \mathcal{I} \}.
\end{array}
\]
For a vector $\vect{\phi}$ of polynomials, 
$\{\vect{\phi}\}$ denotes the polynomial set of entries of $\vect{\phi}$.
\rd{Using these additional constraints, we reformulate
\eqref{eq:problem} as}
\begin{equation}\label{eq:problemhatLME}
	\left\{\begin{array}{cl}
		\max\limits_{\vect{x}\in\re^n} 
		& \sum\limits_{i=1}^m a_i \log p_i(\vect{x})\\
		\st
		& \varphi(\vect{x}) = 0 \, (\varphi\in\Phi),\\
		& \psi(\vect{x})\ge 0 \, (\psi\in\Psi).
	\end{array}\right.
\end{equation}
\rd{Denote the feasible set of \eqref{eq:problemhatLME} by}
\[
K_1 \coloneqq \{\vect{x}\in\re^n: \varphi(\vect{x})=0\,(\varphi\in \Phi),\, 
\psi(\vect{x})\ge 0\,(\psi\in \Psi)\},
\]
\rd{By construction, $K_1\subseteq K$, and the inclusion may be proper.
Under the stated constraint qualification, the global maximizer
$\vect{x}^*$ belongs to $K_1$. Therefore,
\eqref{eq:problemhatLME} has the same optimal value as
\eqref{eq:problem}, although its feasible set may be smaller.}
The problem \eqref{eq:problemhatLME} is also 
a log-polynomial optimization problem.
\rd{Its convex relaxation, analogous to \eqref{eq:relax_meas}, is}
\begin{equation}  \label{eq:relax_meas_lme}
	\left\{\begin{array}{rl}
		f_{\operatorname{lme}} \coloneqq \max\limits_{\vect{z}} 
		&  \sum\limits_{i=1}^m a_i \log \langle p_i,\vect{z}\rangle   \\
		\st & z_0 = 1,\, \vect{z}  \in  \mathscr{R}_d(K_1) .
	\end{array}\right.
\end{equation}
Since $K_1\subseteq K$, we must have 
$f_{\operatorname{lme}} \le f_{\operatorname{mom}}$,
where $f_{\operatorname{mom}}$ is the optimal value of the 
standard convex relaxation \eqref{eq:relax_meas}.
\rd{A hierarchy of semidefinite moment relaxations of
\eqref{eq:relax_meas_lme}, analogous to \eqref{eq:relax_log}, is
constructed as follows. Let}
\[
d_1
\coloneqq
\max\left\{
d,\,
\max_{q\in\Phi\cup\Psi}
\left\lceil\frac{\deg(q)}{2}\right\rceil
\right\}.
\]
\rd{For each $k\geq d_1$, the $k$-th order moment relaxation of
\eqref{eq:relax_meas_lme} is}
\begin{equation}\label{eq:relax_lme}
	\left\{\begin{array}{rl}
		f_{\operatorname{lme},k} \coloneqq \max\limits_{\vect{y}} 
		& \sum\limits_{i=1}^m a_i \log \langle p_i, \vect{y} \rangle \\
		\st 
		& y_0 = 1,\, \vect{y} \in \mathbb{R}^{\mathbb{N}_{2k}^n},\, 
		\mat{M}_k [\vect{y}] \succeq 0,\\
		& \mat{L}_{\varphi}^{(k)}[\vect{y}] = 0\,(\varphi\in\Phi),\\ 
		& \mat{L}_{\psi}^{(k)}[\vect{y}] \succeq 0 \,(\psi  \in \Psi).
	\end{array}\right.
\end{equation}
Since \eqref{eq:relax_lme} includes the constraints of
\eqref{eq:relax_log} together with additional LME constraints, its
feasible set is a subset of the feasible set of
\eqref{eq:relax_log}.
Thus, one has $f_{\operatorname{lme},k}\le f_{\operatorname{mom},k}$ 
for all $k\ge d_1$.


\rd{The construction extends to more general feasible sets whenever
rational LMEs with denominators strictly positive on $K$ are
available.}
Suppose the linear independence constraint qualification (LICQ)
holds for all $\vect{x}\in K$ and all $p_i$ are positive on $K$.
Then there exist rational LMEs with denominators strictly positive on $K$, 
by \cite[Proposition~3.6]{nie2023convex}. 
In summary, we have the following result.
\begin{proposition}\label{prop:rel}
	Suppose there exist rational LMEs $\vect{\tau} = (\tau_1,\ldots, \tau_l)$ that
	satisfy \eqref{eq:lme} with denominators strictly positive on $K$.
	\rd{Then, for every $k\geq d_1$, the convex relaxation
	\eqref{eq:relax_meas_lme} and the $k$-th order SDP relaxation
	\eqref{eq:relax_lme} are relaxations of \eqref{eq:problem}, and their
	optimal values satisfy}
	\[
	\rd{\begin{aligned}
		f_{\max}
		&\leq f_{\operatorname{lme}}
		\leq f_{\operatorname{mom}}
		\leq f_{\operatorname{mom},k},\\
		f_{\max}
		&\leq f_{\operatorname{lme}}
		\leq f_{\operatorname{lme},k}
		\leq f_{\operatorname{mom},k}.
	\end{aligned}}
	\]
\end{proposition}
\rd{It may happen that
\[
f_{\operatorname{lme}}=f_{\max}<f_{\operatorname{mom}}.
\]
\rd{In this case, the standard moment hierarchy
\eqref{eq:relax_log} cannot converge to $f_{\max}$,}
because
\[
f_{\operatorname{mom},k}
\geq f_{\operatorname{mom}}
>f_{\max}
\qquad\text{for every }k.
\]
\rd{Moreover, the LME convex relaxation}
\eqref{eq:relax_meas_lme} has no convexification gap.
If
$\operatorname{Ideal}[\Phi]+\operatorname{QM}[\Psi]$
is Archimedean, Theorem~\ref{thm:asymp}, applied with $K$ replaced
by $K_1$, gives
\[
f_{\operatorname{lme},k}
\longrightarrow f_{\operatorname{lme}}.
\]
Thus, in the case above, the LME hierarchy
\eqref{eq:relax_lme} converges to $f_{\max}$.
}

\rd{The flat-truncation and rank-one extraction results in
Theorem~\ref{tm:parallel_all} and Corollary~\ref{cor:rank1} also carry
over to the LME hierarchy \eqref{eq:relax_lme}, with the representing
measure supported on $K_1$.}

\rd{The concavity-based result in
Theorem~\ref{thm:concave} carries over only when $K_1$ is convex and
the polynomials $p_i$ are concave on $K_1$. The additional
polynomial optimality conditions defining $K_1$ need not preserve
convexity, so these assumptions cannot be expected in general.}

We next compare \eqref{eq:relax_log} and \eqref{eq:relax_lme}
numerically.


\begin{example}\label{ex:ABO}
	Consider the log-polynomial optimization problem derived from the 
	ABO blood group system \cite{EMBook}:
	\begin{equation}\label{eq:ABO}
		\left\{\begin{array}{cl}
			\max\limits_{\vect{x}\in\re^3} 
			& 182 \log (x_1^2+2x_1x_3)+ 60 \log (x_2^2+2x_2x_3)\\
			&\qquad + 17 \log(2x_1x_2) + 176 \log (x_3^2)\\
			\st & 1-\mathbf{1}^T\vect{x} \ge 0,\,\, \vect{x} \geq \mathbf{0},
		\end{array}\right.
	\end{equation}
	where $\mathbf{1}\in\re^3$ is the vector of all ones and 
	$\mathbf{0}\in\re^3$ is the vector of all zeros.
	This optimization problem has simplex constraints, 
	and its LMEs are explicitly given in Example~\ref{ex:lmes} (ii). 
	We compare the bounds and optimizer-extraction behavior of the
	standard moment relaxation \eqref{eq:relax_log} and the tighter LME
	moment relaxation \eqref{eq:relax_lme}.
	The moment relaxations were constructed using
	\texttt{GloptiPoly}~\cite{GloptiPoly}, modeled in
	\texttt{YALMIP}~\cite{Yalmip}, and solved with the semidefinite
	programming (SDP) solver \texttt{MOSEK}~\cite{mosek}.
{\color{red}
For $k=2,\ldots,6$, the standard moment relaxation
\eqref{eq:relax_log} gives
$f_{\operatorname{mom},k}=-491.8158$, whereas the LME moment
relaxation \eqref{eq:relax_lme} gives
$f_{\operatorname{lme},k}=-492.5353$, up to the displayed numerical
precision. No flat truncation is detected for the standard moment
relaxation through order $k=6$. For the LME moment relaxation,
 the lowest relaxation order is $d_1=2$, and its order-$2$ moment
optimizer satisfies the flat truncation condition
\eqref{eq:flat_truncation} with $t=2$ and $r=1$. Consequently,
\[
f_{\max}=f_{\operatorname{lme},2}=-492.5353,
\qquad
\vect{x}^*=(0.2644,0.0932,0.6424).
\]
}
\end{example}

{\color{red}
Choosing between \eqref{eq:relax_log} and
\eqref{eq:relax_lme} involves a tradeoff between bound quality and
computational cost. By definition, $d_1\geq d$. At every common
relaxation order $k\geq d_1$, the feasible set of the LME moment
relaxation \eqref{eq:relax_lme} is contained in that of the standard
moment relaxation \eqref{eq:relax_log}, and hence
\[
f_{\operatorname{lme},k}
\leq
f_{\operatorname{mom},k}.
\]
Thus, \eqref{eq:relax_lme} provides an upper bound at least as tight as
\eqref{eq:relax_log} at the same order. This comparison concerns only
bound quality; the additional LME constraints and the potentially
higher minimum order $d_1$ can substantially increase computational
cost.

For simple feasible sets such as boxes, simplices, and balls, explicit
LMEs are available, and the LME hierarchy \eqref{eq:relax_lme}
may provide substantially tighter bounds when $m$ is small and
$d_1$ remains close to $d$. As $m$
increases, however, the degrees of the products $p^{\mathbf{1}}$ and
$f_i$ can increase, causing both $d_1-d$ and the size of the initial
LME relaxation to grow. In this regime, the standard moment hierarchy
\eqref{eq:relax_log} may be computationally preferable, although it can
yield a substantially looser upper bound. Section~\ref{sec:scaling}
studies this tradeoff by varying $m$, $n$, and the coefficients
$a_i$. It compares the standard moment hierarchy with the LME,
direct-product Moment--SOS, and sums-of-AM/GM-exponentials (SAGE)
relaxations whenever the corresponding formulations are computationally
applicable.
}



\section{Numerical experiments and applications}
\label{sec:exp}


In this section, we present explicit numerical examples 
and applications of log-polynomial optimization.
Unless otherwise stated, the \textsc{Matlab} moment relaxations in
this section were constructed using
\texttt{GloptiPoly}~\cite{GloptiPoly} and
\texttt{YALMIP}~\cite{Yalmip} and solved with
\texttt{MOSEK} 10.0.40~\cite{mosek}.
The computations were implemented in \textsc{Matlab} R2025b on a
laptop equipped with an Intel Core i7-1270P CPU running at
2.20 GHz and 32 GB of RAM.
Unless otherwise stated, numerical ranks are computed with
tolerance $10^{-4}$.
To enhance readability, the floating-point results are reported to
four decimal places.


\subsection{Explicit numerical examples}

\begin{example}\label{ex:88}
	Consider the log-polynomial optimization problem:
	\begin{equation*}
		\left\{\begin{array}{cl}
			\max\limits_{\vect{x} \in \mathbb{R}^5} & 
			\frac{30}{218} \log p_1(\vect{x}) + \frac{97}{218} \log p_2(\vect{x}) + \frac{91}{218} \log p_3(\vect{x}) \\
			\st & \| \vect{x} \| \leq 1,
		\end{array}\right.
	\end{equation*}
	where
	\begin{equation*}
		\begin{array}{rl}
			p_1(\vect{x})  & = (x_3+x_4)^2 + \left(1+2(x_1+x_5)+3(x_2+x_4) \right)^2 + 0.01, \\
			p_2(\vect{x})  & = (x_2+x_4-x_5)^2 + (2x_2+3x_5)^2 + 0.02,  \\
			p_3(\vect{x})  & = (x_1+x_3+x_4)^2 + (x_1-x_3+x_4)^2 + 0.03.
		\end{array}
	\end{equation*} 
	It is clear that the objective function is well defined 
	over the feasible set. 
	\rd{The relaxation \eqref{eq:relax_log} at $k=2$ was solved in
	$1.04$ seconds.}
	The corresponding moment optimizer satisfies the flat truncation 
	condition \reff{eq:flat_truncation} at $t=2$ with $r=1$.
	By Corollary \ref{cor:rank1}, the moment relaxation is tight.
	We obtain the global optimal value
	$f_{\max} = f_{\operatorname{mom},2} = 1.6292$, 
	and the global maximizer
	\begin{equation*}
		\vect{x}^* = (0.4717, 0.4518, 0.0036, 0.5276, 0.5432).
	\end{equation*} 
\end{example}


\begin{example}\label{ex:44}
	Consider the log-polynomial optimization problem:
	\begin{equation}\label{eq:44}
		\left\{\begin{array}{cl}
			\max\limits_{\vect{x} \in \mathbb{R}^3} & \sum\limits_{i=1}^{6} a_i \log  p_i(\vect{x})  \\
			\st & 1 - \mathbf{1}^T \vect{x} \geq 0,\,\, \vect{x} \geq \mathbf{0},
		\end{array}\right.
	\end{equation}
	where $\vect{a} = (a_1,\ldots, a_6)$ is a given coefficient vector and
	\begin{equation*}
		\begin{array}{lll}
			p_1(\vect{x}) = x_1^3 + 3x_1^2x_2 + 3x_1^2 x_3, 
			& p_2(\vect{x}) = 3x_1x_2^2 + 6x_1x_2x_3 , 
			& p_3(\vect{x}) = 3x_1x_3^2,\\
			p_4(\vect{x}) = x_2^3 + 3x_2^2x_3, 
			& p_5(\vect{x}) = 3x_2x_3^2 , 
			& p_6(\vect{x}) = x_3^3 .
		\end{array}
	\end{equation*}
	The problem \eqref{eq:44} has simplex constraints, 
	and its LMEs are explicitly presented in Example~\ref{ex:lmes} (ii).
    In the computations for this example, we add the redundant ball
constraint $\|\vect{x}\|^2\leq 1$.
	
	
	(i) Consider the coefficient vector
	\[
	\vect{a}=
	( 0.0968,\,    0.1419,\,    0.2194,\,    0.0839,\,    0.2839 ,\,   0.1742).
	\]
	
{\color{red}
	For the standard moment relaxation \eqref{eq:relax_log}, the lowest
	relaxation order is $d=2$. At relaxation order $k=4$, the relaxation
	was solved in $3.49$ seconds and gave
	$f_{\operatorname{mom},4}=-1.7181$. The associated moment optimizer
	satisfies the flat truncation condition
	\eqref{eq:flat_truncation} with $t=3$ and $r=2$, admitting the
	decomposition \eqref{eq:ft_decomp} with
	\[
	\begin{array}{ll}
	\vect{v}_1^*=(0.2092,0.2001,0.5906),
	& f(\vect{v}_1^*)=-1.7248,\\
	\vect{v}_2^*=(0.0001,0.3479,0.6519),
	& f(\vect{v}_2^*)=-5.2817.
	\end{array}
	\]
	By Theorem~\ref{tm:parallel_all},
	$f_{\operatorname{mom}}=f_{\operatorname{mom},4}$. However, this
	equality does not imply that
	$f_{\operatorname{mom}}=f_{\max}$. Since $\vect{v}_1^*$ is feasible
	for \eqref{eq:44}, we have
	\[
	0\leq f_{\operatorname{mom}}-f_{\max}
	\leq f_{\operatorname{mom},4}-f(\vect{v}_1^*)
	\leq 0.0067.
	\]

	For the LME moment relaxation \eqref{eq:relax_lme}, the lowest
	relaxation order is $d_1=3$. Flat truncation is not detected at
	$k=3$. At order $k=4$, the relaxation was solved in $4.25$ seconds
	and gave $f_{\operatorname{lme},4}=-1.7193$. Its moment optimizer
	satisfies \eqref{eq:flat_truncation} with $t=4$ and $r=2$, and the
	extracted points are
	\[
	\begin{array}{ll}
	\vect{w}_1^*=(0.1872,0.2161,0.5967),
	& f(\vect{w}_1^*)=-1.7194,\\
	\vect{w}_2^*=(0.9982,0.0013,0.0005),
	& f(\vect{w}_2^*)=-16.1046.
	\end{array}
	\]
	The LME analogue of Theorem~\ref{tm:parallel_all} gives
	$f_{\operatorname{lme}}=f_{\operatorname{lme},4}$. The two extracted
	points do not have the same polynomial image, so flat truncation does
	not certify that $f_{\operatorname{lme}}=f_{\max}$. Nevertheless,
	$\vect{w}_1^*$ is feasible, and hence
	\[
	0\leq f_{\max}-f(\vect{w}_1^*)
	\leq f_{\operatorname{lme},4}-f(\vect{w}_1^*)
	\leq 1.4\times10^{-4}.
	\]
	Thus, the LME relaxation provides a substantially tighter upper
	bound, although it does not eliminate the convexification gap in
	this instance.
}
	
	
	
\clearpage
{\color{red}
	(ii) We use five fixed coefficient vectors.
	Table~\ref{tab:ex:44} reports the computational results for these
	instances. For each relaxation, $k_{\operatorname{mom}}$ or
	$k_{\operatorname{lme}}$ denotes the lowest order at which the flat
	truncation condition \eqref{eq:flat_truncation} is detected, and
	$r_{\operatorname{mom}}$ or $r_{\operatorname{lme}}$ denotes
	$\operatorname{rank}M_t[\vect{y}^*]$ at that flat-truncation order.
}

	\begin{table}[!h]
		\centering
		\caption{Computational results for Example~\ref{ex:44} (ii).}
        \label{tab:ex:44}
\begin{tabular}{|c|c|c|c|c|c|c|}
\hline
& \multicolumn{3}{c|}{Relaxation~\eqref{eq:relax_log}}
& \multicolumn{3}{c|}{Relaxation~\eqref{eq:relax_lme}}\\
\hline
Instance
& $k_{\operatorname{mom}}$
& $r_{\operatorname{mom}}$
& $f_{\operatorname{mom},k_{\operatorname{mom}}}$
& $k_{\operatorname{lme}}$
& $r_{\operatorname{lme}}$
& $f_{\operatorname{lme},k_{\operatorname{lme}}}$\\
\hline
\#1 & 3 & 2  & -4.6824 & 3 & 1 & -4.6968\\
\hline
\#2 & 3 & 2  & -3.9521 & 3 & 1 & -3.9658\\
\hline
\#3 & 8 & 12 & -2.7204 & 4 & 2 & -2.8517\\
\hline
\#4 & 3 & 2  & -6.6434 & 3 & 1 & -6.7018\\
\hline
\#5 & 5 & 3  & -7.0067 & 3 & 1 & -7.0672\\
\hline
\end{tabular}
	\end{table}

{\color{red}
	For instances $\#1$, $\#2$, $\#4$, and $\#5$, the LME moment
	optimizer has rank one. Corollary~\ref{cor:rank1}, as extended to the
	LME hierarchy, therefore certifies
	$f_{\max}=f_{\operatorname{lme},3}$. For instance $\#3$, the LME
	moment optimizer has rank two at $k=4$. Flat truncation certifies
	$f_{\operatorname{lme}}=f_{\operatorname{lme},4}$, but the extracted
	points do not have the same polynomial image, so global exactness is
	not certified for this instance.
	For instances $\#1$, $\#2$, $\#4$, and $\#5$, flat truncation for the
	standard relaxation also gives
	$f_{\operatorname{mom}}=f_{\operatorname{mom},k_{\operatorname{mom}}}$.
	Thus, the differences between the standard and LME values in the table,
	namely $0.0144$, $0.0137$, $0.0584$, and $0.0605$, respectively, are
	the corresponding convexification gaps
	$f_{\operatorname{mom}}-f_{\max}$.

	At the reported relaxation orders, the times required to solve
	\eqref{eq:relax_log} on instances $\#1$--$\#5$ were, respectively,
	$0.97$, $2.22$, $7.40$, $0.41$, and $0.45$ seconds. The corresponding
	times required to solve \eqref{eq:relax_lme} were
	$1.60$, $0.64$, $0.76$, $0.51$, and $0.71$ seconds. These values
	exclude preliminary solves at smaller orders used to identify the
	lowest order satisfying \eqref{eq:flat_truncation}. Because the two
	relaxations are not always solved at the same order, these runtimes
	are computational records rather than a same-order speed comparison.
}
\end{example}

{\color{red}
\begin{example}[An LME relaxation with a rank-two optimal moment matrix]
\label{ex:lme_rank_two}
Consider the symmetric log-polynomial optimization problem
\[
\left\{
\begin{array}{cl}
\displaystyle\max_{x\in\mathbb{R}}
& \log(2+x^2)+\log(3+x^2)\\
\st & -1\leq x\leq 1.
\end{array}
\right.
\]
Its global optimal value and maximizers are
\[
f_{\max}=\log(12),
\qquad x^*\in\{-1,1\}.
\]

To construct the LME moment relaxation \eqref{eq:relax_lme}, let
\[
p_1(x)=2+x^2,\qquad p_2(x)=3+x^2,
\]
Since
\[
f'(x)=\frac{p_1'(x)}{p_1(x)}+\frac{p_2'(x)}{p_2(x)},
\]
multiplying by $p_1(x)p_2(x)$ gives the polynomialized derivative
\[
g(x)
=p_2(x)p_1'(x)+p_1(x)p_2'(x)
=2x(5+2x^2).
\]
For the box constraints, the polynomialized LME multipliers
$\widehat{\tau}_i=p_1p_2\tau_i$ are
\[
\widehat{\tau}_1(x)=\frac{g(x)(x-1)}{2},
\qquad
\widehat{\tau}_2(x)=\frac{g(x)(x+1)}{2}.
\]
The complementarity constraints have degree five, so the lowest
relaxation order of \eqref{eq:relax_lme} is $d_1=3$.

For the standard moment relaxation \eqref{eq:relax_log}, the order-$1$
relaxation is unbounded under this representation of the box constraints.
At $k=2$, the relaxation value agrees numerically with $\log(12)$, but
the flat truncation condition \eqref{eq:flat_truncation} is not satisfied.
The condition first holds at $k=3$, with $t=2$ and rank two. At $k=3$,
the moment optimizer of the LME moment relaxation \eqref{eq:relax_lme} does not satisfy
\eqref{eq:flat_truncation}: numerically,
$\operatorname{rank}M_3=2$ whereas $\operatorname{rank}M_0=1$.
At $k=4$, however,
$\operatorname{rank}M_4=\operatorname{rank}M_1=2$, so the flat
truncation condition holds with $t=4$.
The results are reported in Table~\ref{tab:lme_rank_two}.
\begin{table}[htb]
\centering
\caption{Standard and LME moment relaxations for
Example~\ref{ex:lme_rank_two} at their first relaxation orders
satisfying the flat truncation condition.}
\label{tab:lme_rank_two}
\begin{tabular}{|c|c|c|c|c|c|}
\hline
Relaxation
& Value
& Flat order $t$
& Rank
& Extracted atoms
& Time (seconds)\\
\hline
Standard ($k=3$)
& 2.4849
& 2
& 2
& $-0.9999,\ 0.9999$
& 1.66\\
\hline
LME ($k=4$)
& 2.4849
& 4
& 2
& $-0.9999,\ 0.9999$
& 1.28\\
\hline
\end{tabular}
\end{table}

The differences $|f_{\operatorname{mom},3}-\log(12)|$ and
$|f_{\operatorname{lme},4}-\log(12)|$ are of order $10^{-9}$ and are
attributable to numerical solver tolerances.
For the LME moment relaxation \eqref{eq:relax_lme},
$\operatorname{rank}M_4=2$ at each of the tested tolerances
$10^{-6}$, $10^{-5}$, $10^{-4}$, and $10^{-3}$.

For both relaxations, the extracted atoms have the same polynomial
image because
\[
\mathcal{P}(-1)=\mathcal{P}(1)=(3,4).
\]
Thus, Theorem~\ref{tm:parallel_all}, together with its LME counterpart
described above, explains why the rank-two flat truncations certify
exactness and recover both global maximizers. In particular, an optimal
solution of the LME moment relaxation \eqref{eq:relax_lme} need not
have a rank-one moment matrix.
The standard moment relaxation \eqref{eq:relax_log} already attains
the exact value at $k=2$ and provides a flat-truncation certificate at
$k=3$. Thus, the LME constraints do not improve the bound; instead,
this example provides a direct comparison of the ranks and extracted
solutions produced by the two relaxations.
\end{example}
}


\begin{example}\label{ex:92}
	Consider the polynomial optimization problem:
	\[
	\left\{\begin{array}{cl}
		\max\limits_{\vect{x}\in\re^5} 
		& \theta(\vect{x}) = \big((x_3+x_4+x_5)^2-x_1x_2\big)^{20}\big(8-(x_5-x_2)^2\big)^{25}\\
		\st & \|\vect{x}\|_{2}\le 2,\, (x_3+x_4+x_5)^2-x_1x_2\ge 0,\\
		&  8-(x_5-x_2)^2\ge 0.
	\end{array}
	\right.
	\]
	Since $\deg(\theta) = \rd{90}$, \rd{a direct Moment--SOS relaxation
	starts at order $45$ and is computationally prohibitive.}
	Instead, we can solve the problem by reformulating it into 
	an equivalent log-polynomial optimization problem.
	We consider the following two reformulations:
	\begin{equation}\label{eq:ex92}
		\left\{\begin{array}{cl}
			\max\limits_{\vect{x} \in \mathbb{R}^5} &  20\log\left((x_3+x_4+x_5)^2-x_1x_2\right)+25\log\left(8-(x_5-x_2)^2\right)\\
			\st & \|\vect{x}\|_{2}\le 2,\, (x_3+x_4+x_5)^2-x_1x_2\ge 0,\\
			&  8-(x_5-x_2)^2\ge 0;
		\end{array}\right.
	\end{equation}
	\begin{equation}
		\label{eq:ex92_2}
		\left\{\begin{array}{cl}
			\max\limits_{\vect{x} \in \mathbb{R}^5} &  20\log\left((x_3+x_4+x_5)^2-x_1x_2\right)\\
			& \qquad +25\log\left(2\sqrt{2}-(x_5-x_2)\right)
			+25\log\left(2\sqrt{2}+(x_5-x_2)\right)\\
			\st & \|\vect{x}\|_{2}\le 2,\, (x_3+x_4+x_5)^2-x_1x_2\ge 0,\\
			&  8-(x_5-x_2)^2\ge 0.
		\end{array}\right.
	\end{equation}
	We consider the standard moment relaxation \eqref{eq:relax_log} of 
	\eqref{eq:ex92}--\eqref{eq:ex92_2},
	since these log-polynomial optimization problems have
	relatively complex constraints.
	
	
	(i) For the problem \eqref{eq:ex92}, \rd{the standard moment relaxation
	\eqref{eq:relax_log} was solved in $3.37$ seconds} at relaxation order
	$k=2$, and $f_{\operatorname{mom},2}=99.7252$.
	The corresponding moment optimizer satisfies the flat truncation 
	condition \eqref{eq:flat_truncation} at $t=2$.
	It admits the decomposition \eqref{eq:ft_decomp} with 
	\[\begin{array}{ll}
		\vect{v}_1^* = ( 0.0673,
		-0.3654,
		-1.2400,
		-1.2400,
		-0.8870),\\
		\vect{v}_2^* =  (-0.0673,
		0.3654,
		1.2400,
		1.2400,
		0.8870).
	\end{array}
	\]
	Denote the polynomials $p_1(\vect{x}) = (x_3+x_4+x_5)^2-x_1x_2$
	and $p_2(\vect{x}) = 8-(x_5-x_2)^2$.
	\rd{Since $\vect{v}_2^*=-\vect{v}_1^*$ and
	$p_j(-\vect{x})=p_j(\vect{x})$ for $j=1,2$, we have
	$\mathcal{P}(\vect{v}_1^*)=\mathcal{P}(\vect{v}_2^*)$.}
	Therefore, by Theorem~\ref{tm:parallel_all}, we have
	\[
	f_{\max} = f_{\operatorname{mom},2} = 99.7252
	\]
	and both $\vect{v}_1^*,\vect{v}_2^*$ are global maximizers of the
	original polynomial optimization problem.
	\rd{Consequently, the original polynomial optimization problem has
	optimal value $\exp(99.7252)$.}
	
	
	(ii) For the problem \eqref{eq:ex92_2},
	\rd{the standard moment relaxation \eqref{eq:relax_log} was solved in
	$4.35$ seconds} at relaxation order $k=2$, and
	$f_{\operatorname{mom},2}=101.6842$.
	The corresponding moment optimizer satisfies the flat truncation 
	condition \eqref{eq:flat_truncation} at $t=2$.
	It admits the decomposition \eqref{eq:ft_decomp} with 
	\[\begin{array}{ll}
		\vect{v}_1^* = \rd{(0.0000,
		0.0001,
		-1.1553,
		-1.1553,
		-1.1554)},\\
		\vect{v}_2^* = \rd{(0.0000,
		-0.0001,
		1.1553,
		1.1553,
		1.1554)}.
	\end{array}
	\]
	Denote the polynomials $p_1(\vect{x}) = (x_3+x_4+x_5)^2-x_1x_2$, 
	$p_2(\vect{x}) = 2\sqrt{2}-(x_5-x_2)$ and $p_3(\vect{x}) = 2\sqrt{2}+(x_5-x_2)$.
	\rd{Since $p_2(\vect{v}_1^*)\neq p_2(\vect{v}_2^*)$, the sufficient
	condition in Theorem~\ref{tm:parallel_all} is not satisfied.
	Moreover, \eqref{eq:ex92} and \eqref{eq:ex92_2} have the same
	optimal value $f_{\max}=99.7252$, whereas
	$f_{\operatorname{mom},2}=101.6842$ for \eqref{eq:ex92_2}. Hence,
	\[
	f_{\operatorname{mom},2}-f_{\max}=1.9590,
	\]
	so the moment relaxation of \eqref{eq:ex92_2} is not tight.}
\end{example}









\begin{example}\label{ex:dim10}
	Consider the following log-polynomial optimization problem:
	\begin{equation*}
		\left\{\begin{array}{cl}
			\max\limits_{\vect{x} \in \mathbb{R}^{10}} & \sum\limits_{i=1}^{5} a_i \log p_i(\vect{x}) \\
			\st & p_i(\vect{x}) - i \geq 0, \quad (i \in [5]),\\
			& 6 - (x_6-x_7-x_8)^2 - x_9^2 \geq 0, \\
			& 7 - (x_1-x_{10})^2 - (x_2-x_9)^2 \geq 0,\\
			& (x_1 + x_2 + x_3)^2 + (x_4+x_5)^2 - 8 \geq 0,\\ 
			& (x_2+x_4-x_6+x_8-x_{10})^2 - 9 \geq 0,
		\end{array}\right.
	\end{equation*}
	where each $p_i(\vect{x})$ takes the form
	\[
	p_i(\vect{x})\coloneqq \alpha_i - \big(\beta_i(x_{2i-1} + b_i)^2 
	+ \gamma_i(x_{2i}+c_i)^2 + \delta_i(x_{2i-1}+d_i)(x_{2i}+e_i)\big).
	\]
	The weights $a_i$ and the parameters of $p_i$ are reported in
	Table~\ref{tab:dim10}.
	\begin{table}[htb]
		\centering
		\caption{Parameters of Example~\ref{ex:dim10}}
		\label{tab:dim10}
		\begin{tabular}{|c|c|c|c|c|c|c|c|c|c|}
			\hline
			$i$ & $a_i$ & $\alpha_i$ & $\beta_i$ & $\gamma_i$ & $\delta_i$  & $b_i$ & $c_i$ & $d_i$ & $e_i$\\ \hline
			1 & 3 & 10 & 2 & 3 & -3  & 0.5 & 0.5 & -0.5 & 0.5\\ \hline
			2 & 2 & 12 & 3 & 2 &  -1  & -0.6 & 0.1 & -0.6 & -0.1\\ \hline
			3 & 1 & 11 & 4 & 1 & 0  & 0.7 & -0.2 & 0 & 0 \\ \hline
			4 & 1 & 15 & 2 & 5 & -2 & -0.8 & 0.3 & -0.8 & 0.3 \\ \hline
			5 & 5 & 13 & -3 & 2 & 0  & -0.9 & 0.4 & 0 & 0 \\ \hline
		\end{tabular}
	\end{table}
	By checking the negative semidefiniteness of their Hessian matrices, 
	one can easily verify that $p_1, \ldots, p_4$ are concave, whereas $p_5$ is not.
	\rd{The standard moment relaxation \eqref{eq:relax_log} at $k=2$
	was solved in $3.73$ seconds.}
	The flat truncation condition \eqref{eq:flat_truncation} holds for
	$t=2$ and $r=1$.
	By Corollary \ref{cor:rank1}, 
	the moment relaxation is tight.
	We get the global optimal value $f_{\max} = f_{\operatorname{mom},2} =  36.3994,$
	and the optimizer
	\begin{equation*}
		\begin{array}{rl}
			\vect{x}^* = (&-1.4038, -1.5957, 0.4259, -0.4182, -0.7548, \\
			& 0.6579, 0.9798, -0.3904, -2.4485, -0.0622).
		\end{array}
	\end{equation*}
\end{example}


\begin{example} \label{ex:dim12}
	Consider the log-polynomial optimization problem:
	\begin{equation*}
		\left\{\begin{array}{cl}
			\max\limits_{\vect{x} \in \mathbb{R}^{12}} & \sum\limits_{i=1}^{10} \log p_i(\vect{x}) \\
			\st & \rd{p_i(\vect{x})\geq 0,\quad i\in[10],}\\
			& (x_1 + x_2 + x_3 + x_4 + x_5 + x_6)^2 \leq 15,\\
			& (x_7 + x_8 + x_9 + x_{10} + x_{11} + x_{12})^2 \leq 15,\\
			& (x_1 - x_3 + x_5 - x_7 + x_9 - x_{11})^2 \leq 8, \\
			& (x_2 - x_4 + x_6 - x_8 + x_{10} - x_{12})^2 \leq 8, \\
			& (x_1 + x_{12})^2 + (x_2 + x_{11})^2 \leq 9,\\
			& (x_3 - x_{10})^2 + (x_4 - x_9)^2 + (x_5-x_8)^2 \leq 9, \\
			& \sum\limits_{i=1}^6 x_{2i}^2  + 2(x_1^2 + x_3^2 + x_7^2 + x_9^2) + 3(x_5^2 + x_{11}^2) \leq 20,\\
			& \sum\limits_{i=1}^6 (x_{2i-1} - x_{2i}) \leq 5, \, (x_1 + x_6 + x_{12})^2 \leq 4, \\
			& x_1 - x_{12} \leq 3, \, \mathbf{1}^T \vect{x} \geq 0,
		\end{array}\right.
	\end{equation*}
	In the above, for each $i = 1\ddd 10$,
	\begin{equation*}
		p_i(\vect{x}) \coloneqq \alpha_i - \beta_i \big(h_i(\vect{x})\big)^4,
	\end{equation*}
	where for $i=1,\ldots, 5$,
	\[
	h_i(\vect{x}) \coloneqq (x_{2i}+\xi_i) + (x_{2i+1}+\zeta_i),
	\]
	and for $i=6,\ldots, 10$ ($j_1\bmod j_2$ denotes the remainder of $j_1$ modulo $j_2$)
	\[ \begin{array}{ll}
		h_i(\vect{x}) \coloneqq \;& (-1)^{i-4}(x_{i-5}+\xi_i) 
		+ (-1)^{i+1} (x_i +\zeta_i) \\
		& + (-1)^{(i+4) \bmod 12} (x_{(i+4) \bmod 12 + 1} + \omega_i).
	\end{array}
	\]
	The parameters in $p_i$ and $h_i$ are reported in
	Table~\ref{tab:dim12}.
	\begin{table}[htb]
		\centering
		\caption{Parameters of Example~\ref{ex:dim12}}
		\label{tab:dim12}
		\begin{tabular}{|c|c|c|c|c|c|c|c|c|c|c|}
			\hline
			$i$ & 1 & 2 & 3 & 4 & 5 & 6 & 7 & 8 & 9 & 10 \\ \hline
			$\alpha_i$ & 20 & 25 & 18 & 22 & 30 & 15 & 28 & 19 & 21 & 26 \\ \hline
			$\beta_i$ & 1.5 & 1.2 & 2.5 & 1.1 & 1.6 & 1.9 & 2.3 & 1.7 & 2.4 & 1.4 \\ \hline
			$\xi_i$ & -0.2 & -1 & 0.4 & 0.4 & -1.1 & -0.8 & -0.3 & -1.7 & 1.2 & 1.7 \\ \hline
			$\zeta_i$ & -1.6 & -0.4 & -1 & 0.8 & -1.5 & -0.7 & 0 & -1 & -1.9 & 0.9 \\ \hline
			$\omega_i$ & - & - & - & - & - & -1.3 & 0.7 & -0.7 & 0.8 & -0.5 \\ \hline
		\end{tabular}
	\end{table}
	\rd{Since each $h_i$ is affine and $\beta_i>0$,
	\[
	-\nabla^2p_i(\vect{x})
	=12\beta_i h_i(\vect{x})^2
	\nabla h_i\nabla h_i^T,
	\]
	which is an SOS matrix. Thus, all $p_i$ are SOS-concave.
	The constraints $p_i(\vect{x})\geq0$ are convex because the $p_i$
	are concave, and the remaining constraints also define convex sets.
	Hence, this is a convex optimization problem.}
	\rd{The standard moment relaxation \eqref{eq:relax_log} at $k=2$
	was solved in $9.18$ seconds.}
	The corresponding moment optimizer satisfies the flat truncation
condition \eqref{eq:flat_truncation} at $t=2$ with $r=1$.
By Corollary~\ref{cor:rank1}, the moment relaxation is tight. 
Thus,
	$
	f_{\max} = f_{\operatorname{mom},2} = 30.2624,
	$
	and the obtained global maximizer is 
	\begin{equation*}
	\begin{array}{rl}
		\vect{x}^* = (& -0.3944, -0.8459, 1.5027, -0.9420, 1.0771, -0.4394,\\
		& 0.4717, -1.0822, 1.0639, 2.0626, 0.6715, 0.6394).
	\end{array}
\end{equation*}
\end{example}

{\color{red}
\subsection{Scaling and comparisons with alternative relaxations}
\label{sec:scaling}

We next examine the quality and computational cost of the standard
moment relaxation \eqref{eq:relax_log} and the LME moment relaxation
\eqref{eq:relax_lme} as the number of logarithmic terms, the number of
variables, and the exponents increase. We consider the simplex
\[
K_n\coloneqq\{\vect{x}\in\mathbb{R}^n:\vect{x}\geq 0,\
\mathbf{1}^T\vect{x}\leq 1\}
\]
and the positive quadratic polynomials
\[
p_i(\vect{x})=0.05+(\beta_i+\vect{u}_i^T\vect{x})^2,\qquad i\in[m].
\]
For each instance, the vector $\vect{u}_i$ is generated from a standard Gaussian
vector and normalized to have unit Euclidean norm, while
$\beta_i$ is sampled uniformly from $[-0.15,0.15]$.
The scaling experiment uses $a_i=1$ for every $i\in[m]$.
All generated objectives are nonconcave
on $K_n$. Indeed,
\[
\nabla^2 f(\vect{0})
=
\sum_{i=1}^m
\frac{2(0.05-\beta_i^2)}
{(0.05+\beta_i^2)^2}
\vect{u}_i\vect{u}_i^T.
\]
Since $|\beta_i|\leq0.15<\sqrt{0.05}$ and each $\vect{u}_i$ has unit
Euclidean norm, this matrix is nonzero and positive semidefinite.
It therefore has positive curvature along at least one feasible
coordinate direction near $\vect{x}=\vect{0}$.

For each instance, we use the sequential quadratic programming
(SQP) algorithm of \texttt{fmincon} in \textsc{Matlab} to locally
maximize the original nonconvex objective from 100 fixed feasible
starting points.
Let $\mathcal{X}_{\mathrm{loc}}$ contain these starting
points and all feasible terminal points returned by the local runs, and
define
\[
f_{\mathrm{feas}}
\coloneqq
\max_{\vect{x}\in\mathcal{X}_{\mathrm{loc}}}f(\vect{x}).
\]
Thus, $f_{\mathrm{feas}}\leq f_{\max}$ is the best feasible value found,
not, in general, a certified global optimal value.
Since the polynomials $p_i$ are quadratic, the lowest admissible
order of the standard moment relaxation \eqref{eq:relax_log} is $k=1$.
We report
\[
\Delta_{\mathrm{mom}}
\coloneqq f_{\mathrm{mom},1}-f_{\mathrm{feas}},
\qquad
R_{\mathrm{mom}}
\coloneqq
\exp(\Delta_{\mathrm{mom}})
=
\frac{\exp(f_{\mathrm{mom},1})}{\exp(f_{\mathrm{feas}})}.
\]
The ratio $R_{\mathrm{mom}}$ measures the moment upper bound relative
to the best feasible value on the original product scale.
When the LME moment relaxation \eqref{eq:relax_lme} is computed, we also report
\[
\Delta_{\mathrm{lme}}
\coloneqq
f_{\mathrm{lme},d_1}-f_{\mathrm{feas}},
\]
where $d_1$ is the lowest relaxation order of \eqref{eq:relax_lme}.
For the instances in which the direct Moment--SOS relaxation of
$\prod_i p_i^{a_i}$ satisfies the flat truncation condition
\eqref{eq:flat_truncation}, its certified global value agrees
numerically with $f_{\mathrm{feas}}$.
Otherwise, $\Delta_{\mathrm{mom}}$ and, when reported,
$\Delta_{\mathrm{lme}}$ are only the observed differences between the
corresponding upper bounds and the best feasible value; they are not
claimed to be exact relaxation gaps.

We first vary $m\in\{2,4,6,8,10,12\}$ with $n=4$, and then vary
$n\in\{2,4,6,8\}$ with $m=4$. Table~\ref{tab:scaling} reports one
fixed instance for each pair $(m,n)$. The column $t_{\mathrm{mom}}$
gives the runtime in seconds for solving the standard moment relaxation
\eqref{eq:relax_log}. Model construction and solution extraction are
excluded from the reported runtime.
The columns $s_{\mathrm{lme}}$ and $s_{\mathrm{prod}}$ give the largest
positive semidefinite block dimensions required by the LME moment
relaxation \eqref{eq:relax_lme} and by the direct-product Moment--SOS
relaxation, respectively.

\begin{table}[htb]
\centering
\caption{Scaling behavior on fixed positive nonconcave quadratic
instances.}
\label{tab:scaling}
\begin{tabular}{|c|c|c|c|c|c|c|}
\hline
axis & size & $\Delta_{\mathrm{mom}}$
& $R_{\mathrm{mom}}$
& $t_{\mathrm{mom}}$ (seconds)& $s_{\mathrm{lme}}$ & $s_{\mathrm{prod}}$\\
\hline
$m$ & 2  & 0.6923 & 1.9984   & 2.8014 & 35   & 15\\
$m$ & 4  & 2.0110 & 7.4704   & 0.7645 & 126  & 70\\
$m$ & 6  & 3.6346 & 37.8848  & 0.5331 & 330 & 210\\
$m$ & 8  & 4.5331 & 93.0479  & 0.5232 & 715 & 495\\
$m$ & 10 & 4.6270 & 102.2067 & 0.4536 & 1365 & 1001\\
$m$ & 12 & 6.2265 & 505.9633 & 0.2402 & 2380 & 1820\\
\hline
$n$ & 2 & 0.1534 & 1.1658  & 0.2747 & 21   & 15\\
$n$ & 4 & 2.0110 & 7.4704  & 0.7645 & 126  & 70\\
$n$ & 6 & 2.3800 & 10.8047 & 0.1940 & 462 & 210\\
$n$ & 8 & 1.8332 & 6.2537  & 0.2780 & 1287 & 495\\
\hline
\end{tabular}
\end{table}

The block dimensions in Table~\ref{tab:scaling} were computed
before solving. The LME moment relaxation \eqref{eq:relax_lme} was not
launched when $s_{\mathrm{lme}}>200$, and the direct-product
Moment--SOS relaxation was not launched when
$s_{\mathrm{prod}}>250$.

The standard moment relaxation \eqref{eq:relax_log} remains inexpensive,
but its observed gap
increases substantially with $m$ and varies with $n$.
For the fixed instances with $m=2$ and $m=4$, the LME moment
relaxation \eqref{eq:relax_lme} is computed at its initial orders
$d_1=3$ and $d_1=5$, respectively.
The corresponding values of $\Delta_{\mathrm{lme}}$ are $0.0054$ and
$0.6106$, and the runtimes are $2.1402$ and $4.2927$ seconds.
At $m=6$, the LME block dimension is already 330, whereas the direct
product Moment--SOS relaxation has block dimension 210.
The direct-product relaxation was solved in $6.7448$ seconds, and
its optimal moment sequence satisfies the flat truncation condition
\eqref{eq:flat_truncation}. This illustrates that
the LME strengthening is not
uniformly preferable to the direct product formulation when the
exponents are small integers.

We next fix $n=4$ and $m=3$ and vary the coefficient vector
$\vect{a}=(a_1,a_2,a_3)$. The standard moment relaxation
\eqref{eq:relax_log} always has PSD block dimension 5, while the LME
moment relaxation \eqref{eq:relax_lme} always has order 4 and PSD block
dimension 70. In contrast, when $\vect{a}$ is a positive integer vector, the
minimum order of the direct product Moment--SOS relaxation is
$\sum_i a_i$.

Finally, we compare with the level-zero conditional SAGE relaxation
of Murray et al.~\cite{MurraySAGE}. A conditional SAGE relaxation is
a convex relative-entropy relaxation that certifies nonnegativity on
a specified convex domain. We use the basic relaxation, called level
zero in the SAGE hierarchy; unlike higher levels, it applies the SAGE
certificate directly, without a hierarchy multiplier. For this comparison we use the substitution
$x_i=\exp(y_i)$, under which the simplex becomes the convex log-domain
$\sum_i\exp(y_i)\leq 1$. The SAGE computations use \texttt{sageopt}
0.6.1 and MOSEK \rd{10.0.40}.
For a positive integer coefficient vector $\vect{a}$, let
$f_{\mathrm{SAGE}}$ be the logarithm of the SAGE upper bound for
$\prod_i p_i^{a_i}$, and define
\[
\Delta_{\mathrm{SAGE}}
\coloneqq f_{\mathrm{SAGE}}-f_{\mathrm{feas}}.
\]

\rd{The results are reported in Table~\ref{tab:exponents}.}
\begin{table}[htb]
\centering
\caption{Bound gaps and runtimes for integer and noninteger
logarithmic coefficients.}
\label{tab:exponents}
\renewcommand{\arraystretch}{1.1}
\setlength{\tabcolsep}{3.5pt}
\begin{tabular}{|c|c|c|c|c|c|c|c|c|}
\hline
& \multicolumn{4}{c|}{Bound gaps}
& \multicolumn{4}{c|}{Runtimes (seconds)}\\
\hline
$\vect{a}$ & $\Delta_{\mathrm{mom}}$
& $\Delta_{\mathrm{lme}}$
& $\Delta_{\mathrm{prod}}$
& $\Delta_{\mathrm{SAGE}}$
& $t_{\mathrm{mom}}$
& $t_{\mathrm{lme}}$
& $t_{\mathrm{prod}}$
& $t_{\mathrm{SAGE}}$\\
\hline
$(1,1,1)$
& 1.9954 & 0.2775 & 0 & 1.4444
& 2.2879 & 2.2569 & 1.0937 & \rd{2.6443}\\
\hline
$(1,2,3)$
& 3.7591 & 0.7545 & 0 & 1.3724
& 0.6818 & 2.5681 & 7.0035 & \rd{697.7900}\\
\hline
$(2,5,10)$
& 9.5145 & 1.3227 & NL & NL
& 0.4946 & 4.0863 & NL & NL\\
\hline
$(5,10,20)$
& 20.4593 & 3.1742 & NL & NL
& 1.5051 & 1.7166 & NL & NL\\
\hline
$(0.5,1.5,2.5)$
& 2.6213 & 0.4834 & NA & NA
& 0.3402 & 1.5147 & NA & NA\\
\hline
$(\sqrt{2},\pi/2,e/2)$
& 2.8041 & 0.3197 & NA & NA
& 0.3492 & 1.6082 & NA & NA\\
\hline
\end{tabular}
\end{table}
\FloatBarrier

All rows of Table~\ref{tab:exponents} use the same fixed polynomial
instance. The quantities $\Delta_{\mathrm{mom}}$,
$\Delta_{\mathrm{lme}}$, $\Delta_{\mathrm{prod}}$, and
$\Delta_{\mathrm{SAGE}}$ are differences between the corresponding
upper bounds and $f_{\mathrm{feas}}$, with the direct-product bound
reported on the logarithmic scale. All runtimes are in seconds. The
reported SAGE runtime is the solver time; construction time is excluded.
The two launched direct-product relaxations satisfy
\eqref{eq:flat_truncation}, so their certified gaps are zero. In
Table~\ref{tab:exponents}, NL means not launched and NA means not
applicable. The two larger integer cases were not launched because
their direct-product orders are 17 and 35, while their estimated SAGE
expansions, containing up to $73{,}815$ and $1{,}150{,}626$ monomials,
respectively, were prohibitively large. For noninteger
coefficients, neither the direct-product formulation nor the SAGE
formulation considered here is applicable.

Table~\ref{tab:exponents} shows that the standard and LME
relaxations \eqref{eq:relax_log} and \eqref{eq:relax_lme} remain
inexpensive as the coefficients increase on this fixed instance,
although their observed gaps increase with the coefficient magnitudes.
For the first two integer vectors, the direct-product Moment--SOS
relaxation is exact, and conditional SAGE gives a
tighter bound than the standard moment relaxation
\eqref{eq:relax_log}, but a looser bound than the LME relaxation
\eqref{eq:relax_lme}. Increasing the coefficient vector from
$(1,1,1)$ to $(1,2,3)$ increases the SAGE runtime from about
\rd{2.6} seconds to about 12 minutes.
The direct-product Moment--SOS and SAGE formulations become
prohibitively large as the integer coefficients increase and are not
applicable to noninteger coefficients in the formulations considered
here. Thus, the standard
moment, LME, direct-product Moment--SOS, and conditional SAGE
formulations have different tradeoffs in bound quality, runtime, and
applicability.
}


\subsection{Applications}
\label{sec:app}


\begin{example}[Paternity analysis \cite{McCulloch87}] 
	\label{ex:offspring}
	Consider a single genetic locus involving a single mother 
	and $n$ possible fathers.
	Suppose there are $g$ possible genotypes.
	Let $x_j$ denote the probability that the $j$th male sires an
	offspring for $j=1,\ldots,n$.
	Our goal is to estimate the probabilities $\vect{x}=(x_1,\ldots, x_n)$ 
	using genotype data from a potentially multiply-parented litter.
	This can be done by solving 
	\begin{equation}\label{eq:offspring}
		\left\{\begin{array}{cl}
			\max\limits_{\vect{x}\in\re^n} & \sum\limits_{i=1}^{g} N_i \log \Big( \sum\limits_{j=1}^n P_{ij}x_j \Big)\\
			\st & \mathbf{1}^T\vect{x} = 1,\,\, \vect{x} \geq \mathbf{0},
		\end{array}\right.
	\end{equation}
	where each $N_i$ denotes the number of offspring in the litter with 
	the $i$-th possible genotype,
	and each $P_{ij}$ is the conditional probability that an offspring
	has the $i$-th genotype given that the $j$-th male is its father.

	Let $\vect{\bar{x}} = (x_1,\ldots, x_{n-1})$.
	With the substitution $x_n = 1-\mathbf{1}_{n-1}^T\vect{\bar{x}}$, 
	we can reformulate \eqref{eq:offspring} into 
	\begin{equation}\label{eq:offspring_updated}
		\left\{\begin{array}{cl}
			\max\limits_{\vect{\bar{x}}\in\re^{n-1}} & \sum\limits_{i=1}^{g} N_i\log\Big(
			P_{in}+\sum\limits_{j=1}^{n-1}(P_{ij}-P_{in})x_j\Big)\\
			\st &  \vect{\bar{x}}\geq \mathbf{0},\\
			& 1-\mathbf{1}_{n-1}^T\vect{\bar{x}}\ge 0.
		\end{array}
		\right.
	\end{equation}

This is a log-polynomial optimization problem with simplex constraints.
Since the logarithm arguments are affine and the feasible set is
a simplex, this is a convex optimization problem that can generally
be solved efficiently by standard convex optimization methods. We
therefore include this example as an application and certification
check, rather than as evidence that the moment relaxation is
computationally preferable to such methods.
It can also be solved using the LME moment relaxation
\eqref{eq:relax_lme}, with the LMEs in
Example~\ref{ex:lmes}~(ii); below, we use this formulation to
illustrate the moment-based global optimality certificate and
solution-extraction procedure.

We use genotype-probability matrices $\mat{P}$ based on the
simulation configurations in \cite{McCulloch87}, together with the
fixed count vectors $\vect{N}$ reported in
Table~\ref{tab:offspring}.
	After extracting an optimizer $\vect{\bar{x}}^*$ of
	\eqref{eq:offspring_updated}, we recover an optimizer of the original
	problem \eqref{eq:offspring} as
	$\vect{x}^*=(\vect{\bar{x}}^*,
	1-\mathbf{1}_{n-1}^T\vect{\bar{x}}^*)$.
	For every reported instance, the computed optimizer of the LME
	moment relaxation \eqref{eq:relax_lme} satisfies the flat truncation
	condition \eqref{eq:flat_truncation}, and the corresponding moment
	matrix has rank one. By the LME analogue of
	Corollary~\ref{cor:rank1}, the LME moment relaxation is tight for the
	original problem.

	The detailed input data and our numerical results are reported 
	in Table~\ref{tab:offspring}.
	In the table, we denote parameters
	$\vect{N} = (N_1,\ldots, N_{g})$ 
	and $\mat{P} = (P_{ij})\in\re^{g\times n}$.
	The value $k$ is the relaxation order used in
	\eqref{eq:relax_lme}, $t$ is the order at which the flat truncation
	condition \eqref{eq:flat_truncation} holds, and
	$r=\rank\mat{M}_t[\vect{y}^*]$.
	The notation $\vect{x}^*$ denotes the computed optimizer of the original log-polynomial
	optimization problem \eqref{eq:offspring}.
	The column `runtime' reports the runtime in seconds for solving
	the $k$th-order LME moment relaxation \eqref{eq:relax_lme}.
	After one unreported warm-up solve, each runtime is one wall-clock
	measurement of the corresponding solve.
	\begin{table}[ht]
	\centering
	\caption{Numerical results for Example \ref{ex:offspring}}
    \label{tab:offspring}
	\setlength{\tabcolsep}{3pt}
	\begin{tabular}{|c|c|c|c|c|c|c|c|c|}\hline
		$n$ & $g$ & $\vect{N}$ & $\mat{P}$ & $k$ & $t$ & $r$ & $\vect{x}^*$ & runtime (seconds)\\ \hline
		%Z
		2 & 2 &$\left[\begin{array}{r}
			77 \\ 23
		\end{array}\right]$ & $\begin{bmatrix}
			0.5 & 1 \\
			0.5 & 0
		\end{bmatrix}$ & 2 & \rd{2} & 1 & $\left[\begin{array}{r}
			0.4600 \\ 0.5400
		\end{array}\right]$ & 0.77\\ \hline
		%S
		2 & 2 &$\left[\begin{array}{r}
			63 \\ 37
		\end{array}\right]$ & $\begin{bmatrix}
			0.5 & 1 \\
			0.5 & 0
		\end{bmatrix}$ & 2 & \rd{2} & 1 & $\left[\begin{array}{r}
			0.7400 \\ 0.2600
		\end{array}\right]$ & 2.22\\ \hline
		%B
		2 & 3 &$\left[\begin{array}{r}
			49 \\ 40 \\ 11
		\end{array}\right]$ & $\begin{bmatrix}
			0.5 & 0.875 \\
			0.25 & 0.125 \\
			0.25 & 0
		\end{bmatrix}$ & 3 & \rd{3} & 1 & $\left[\begin{array}{r}
			0.8813 \\ 0.1187
		\end{array}\right]$ & 0.46\\ \hline
		%T
		2 & 3 &$\left[\begin{array}{c}
			83 \\ 2 \\ 15
		\end{array}\right]$ & $\begin{bmatrix}
			0.5 & 0.25\\
			0.5 & 0.5\\
			0 & 0.25
		\end{bmatrix}$ & 2 & \rd{2} & 1 & $\left[\begin{array}{r}
			0.6939 \\ 0.3061
		\end{array}\right]$ & 1.25\\ \hline
		%Y
		2 & 3 &$\left[\begin{array}{r}
			63 \\ 17 \\ 20
		\end{array}\right]$ & $\begin{bmatrix}
			0.5 & 0.25 \\
			0.5 & 0.5 \\
			0 & 0.25
		\end{bmatrix}$ & 2 & \rd{2} & 1 & $\left[\begin{array}{r}
			0.5181 \\ 0.4819
		\end{array}\right]$ & 0.17\\ \hline
		%W
		2 & 4 &$\left[\begin{array}{c}
			59 \\ 8 \\ 16 \\ 17
		\end{array}\right]$ & $\begin{bmatrix}
			0.25 & 0.5 \\
			0.25 & 0.5 \\
			0.25 & 0 \\
			0.25 & 0
		\end{bmatrix}$ & 3 & \rd{3} & 1 & $\left[\begin{array}{r}
			0.6599 \\ 0.3401
		\end{array}\right]$ & 0.23\\ \hline
		%X
		2 & 5 &$\left[\begin{array}{c}
			7 \\ 9 \\ 4 \\ 33 \\ 47
		\end{array}\right]$ & $\begin{bmatrix}
			0.25 & 0 \\
			0.25 & 0 \\
			0.25 & 0 \\
			0.25 & 0.5 \\
			0 & 0.5
		\end{bmatrix}$ & 4 & \rd{4} & 1 & $\left[\begin{array}{r}
			0.2463 \\ 0.7537
		\end{array}\right]$ & 0.11\\ \hline
		%A
		3 & 3 &$\left[\begin{array}{r}
			39 \\ 38 \\ 23
		\end{array}\right]$ & $\begin{bmatrix}
			0.5 & 0 & 0.875 \\
			0.25 & 0.75 & 0.125 \\
			0.25 & 0.25 & 0
		\end{bmatrix}$ & 3 & \rd{2} & 1 & $\left[\begin{array}{r}
			0.6400 \\ 0.2800 \\ 0.0800
		\end{array}\right]$ & 0.31\\ \hline
		%C
		3 & 4 &$\left[\begin{array}{r}
			29 \\ 21 \\ 88 \\ 62
		\end{array}\right]$ & $\begin{bmatrix}
			0.5 & 0 & 0 \\
			0.25 & 0.75 & 0.25 \\
			0.25 & 0.25 & 0.5 \\
			0 & 0 & 0.25
		\end{bmatrix}$ & \rd{4} & \rd{3} & 1 & $\left[\begin{array}{r}
			0.2240 \\ 0.0000 \\ 0.7760
		\end{array}\right]$ & 0.32\\ \hline
	\end{tabular}
\end{table}
\end{example}


\begin{example}[Latent class models \cite{ChenLCM20}]\label{ex:LCM}
	Latent class models (LCMs) are a type of finite 
	mixture model used to identify hidden, unobserved categorical groups 
	within a population based on observed categorical data.
	Let $\vect{y} = (y_1,\ldots, y_d)$ be a vector of categorical random 
	variables, where each $y_j$ takes values from $c_j$ categories, 
	say, $\{1,\ldots, c_j\}$.
	Let $s\in\{1,\ldots,T\}$ denote the latent class variable.
	Conditional on $s$, the categorical variables
	$y_1,\ldots,y_d$ are assumed to be independent.
	Let $I(\cdot)$ denote the Iverson bracket function
	such that $I(P) = 1$ if $P$ is true and zero otherwise.
	Then the probability mass function of $\vect{y}$ in the $t$-th 
	subgroup is
	\[
	\prod\limits_{j=1}^d\prod\limits_{l=1}^{c_j} \pi_{t,j,l}^{I(y_j=l)}
	\quad \mbox{with}\quad
	\pi_{t,j,l} = \Pr(y_j=l\mid s=t),
	\]
	Set
	$\vect{\pi}=(\pi_{t,j,l})_{t\in[T],\,j\in[d],\,l\in[c_j]}$.
	The overall probability mass function of the latent class model
	is then a weighted sum
	\[
	p(\vect{y}\mid\vect{\eta},\vect{\pi}) = \sum\limits_{t=1}^T
	\Big(\eta_t	\prod\limits_{j=1}^d\prod\limits_{l=1}^{c_j} 
	\pi_{t,j,l}^{I(y_j=l)}\Big),
	\]
	where $\vect{\eta}=(\eta_1,\ldots,\eta_T)$ consists of
	nonnegative mixing weights
	satisfying $\sum_{t=1}^T\eta_t=1$.
	Consider $N$ samples drawn from the LCM distribution, denoted as $\vect{y}^{(1)}, \ldots, \vect{y}^{(N)}$.
	The maximum-likelihood estimation problem for the LCM is then
	formulated as follows:
	\begin{equation}\label{eq:LCM}
		\left\{\begin{array}{cl}
			\max\limits_{\vect{\eta}, \vect{\pi}} & \sum\limits_{i=1}^N \log \left(\sum\limits_{t=1}^T \eta_t	\prod\limits_{j=1}^d\prod\limits_{l=1}^{c_j} \pi_{t,j,l}^{I(y_j^{(i)}=l)}\right)\\
			\st & \vect{\eta}\ge 0,\, \vect{\pi}\ge 0,\, \sum\limits_{t=1}^T\eta_t = 1,\\
			& \sum\limits_{l=1}^{c_j} \pi_{t,j,l} = 1\, ( t\in[T],\, j\in[d]).
		\end{array}
		\right.
	\end{equation}
We consider three instances using the same settings as in \cite{ChenLCM20}.
We call an LCM setting saturated when its model image fills the
corresponding probability simplex, so every empirical distribution can
be represented by the model. Cases (i)--(iii) below are saturated in
this sense. They serve as reproducibility and implementation checks but
do not test genuinely non-saturated behavior.
For reproducibility, we use one fixed simulated dataset of $N=500$
observations for each instance and approximate \eqref{eq:LCM}
by its standard moment relaxation \eqref{eq:relax_log} at the
initial order $k=2$.
For cases (i)--(iii), we compare the resulting moment upper bound
with the log-likelihood at the data-generating parameters reported in
\cite{ChenLCM20}. We then add a fourth, non-saturated binary LCM instance that exhibits a persistent convexification gap.
In the binary examples below, the outcomes $0$ and $1$
correspond to the parameter indices $l=1$ and $l=2$, respectively.
After one unreported warm-up solve, we measured one wall-clock runtime
for each of cases (i)--(iii).
	
	(i) For $d=1, T = 2$, suppose $y_1$ is a binary variable, 
	i.e., $y_1\in\{0,1\}$. Let
	\[\begin{array}{llll}
		\eta_1 = 0.5, &  \eta_2 = 0.5, &  \pi_{1,1,1} = 0.4, &  \pi_{1,1,2} = 0.6,\\
		\pi_{2,1,1} = 0.8, & \pi_{2,1,2} = 0.2.
	\end{array}\] 
    In the fixed dataset, 303 observations have $\vect{y}=0$ and 197 observations have $\vect{y}=1$.
	\rd{The standard moment relaxation \eqref{eq:relax_log} was solved in
	$1.72$ seconds}, and the moment upper bound was
	$f_{\operatorname{mom},2}=-335.2519$.
This result is very close to the log-likelihood $-335.29$ at the
data-generating parameters.
	
	(ii) For $d=1, T=3$, suppose $y_1$ is a binary variable, 
	i.e., $y_1\in\{0,1\}$. Let
	\[
	\begin{array}{lllll}
		\eta_1 = 0.5, & \eta_2 = 0.3, & \eta_3 = 0.2, & \pi_{1,1,1} = 0.4, &  \pi_{1,1,2} = 0.6,\\
		\pi_{2,1,1} = 0.8, & \pi_{2,1,2} = 0.2 & \pi_{3,1,1} = 0.1, & \pi_{3,1,2} = 0.9.
	\end{array}
	\]
    In the fixed dataset, 227 observations have $\vect{y}=0$ and 273 observations have $\vect{y}=1$.
	\rd{The standard moment relaxation \eqref{eq:relax_log} was solved in
	$4.36$ seconds}, and the moment upper bound was
	$f_{\operatorname{mom},2}=-344.45$.
This result is very close to the log-likelihood $-344.49$ at the
data-generating parameters.
	
	
	(iii) For $d=2, T=2$, suppose each $y_j$, $j=1,2$, is a binary variable, 
	i.e., $y_j\in\{0,1\}$. Let
	\[
	\begin{array}{llllll}
		\eta_1 = 0.5, & \eta_2 = 0.5,  & \pi_{1,1,1} = 0.4, &  \pi_{1,1,2} = 0.6,
		& \pi_{2,1,1} = 0.8,\\
		\pi_{2,1,2} = 0.2, & \pi_{1,2,1} = 0.1, & \pi_{1,2,2} = 0.9 & \pi_{2,2,1} = 0.6, & \pi_{2,2,2} = 0.4.
	\end{array}
	\]
    \rd{For the outcomes $(0,0)$, $(0,1)$, $(1,0)$, and $(1,1)$,
    the fixed dataset has counts 131, 174, 49, and 146, respectively.}
	\rd{The standard moment relaxation \eqref{eq:relax_log} was solved in
	$2.73$ seconds}, and the moment upper bound was
	$f_{\operatorname{mom},2}=-652.6718$.
This result is close to \rd{the log-likelihood $-653.1616$ at the
data-generating parameters}.

{\color{red}
\medskip

\noindent
(iv) We next consider a genuinely non-saturated binary LCM with
$d=4$ and $T=2$. The data-generating mixing proportions are
$\eta_1=\eta_2=0.5$, and the conditional success probabilities are
\[
\left(\Pr(y_j=1\mid s=t)\right)_{\substack{t=1,2\\j=1,\ldots,4}}
=
\begin{pmatrix}
0.38 & 0.42 & 0.46 & 0.50\\
0.62 & 0.58 & 0.54 & 0.50
\end{pmatrix}.
\]
We use a fixed sample of $N=500$ observations. In lexicographic order
\[
(0000),(0001),\ldots,(0111),(1000),\ldots,(1111),
\]
the sixteen outcome counts are
\[
\begin{array}{c}
(30,26,39,26,27,19,28,30,\\
\phantom{(}24,25,39,33,45,31,35,43).
\end{array}
\]
This dataset was selected as a diagnostic stress instance from
30 fixed candidate datasets using a 20-start expectation--maximization
(EM) screening run \cite{Dempster1977,EMBook}. The candidates
were ranked first by the
fraction of starts attaining the best terminal value and then by the
separation from the second-best terminal value. Thus, this example is
intended as a stress test rather than as a representative average-case
experiment.
The EM algorithm was implemented directly in \textsc{Matlab}.

The cell-probability polynomials have degree five, so the lowest
relaxation order is $k=3$. The numerical results are reported
in Table~\ref{tab:lcm_nonsaturated}.
\begin{table}[htb]
\centering
\caption{Moment upper bound and multistart EM result for the
non-saturated LCM in Example~\ref{ex:LCM}(iv).}
\label{tab:lcm_nonsaturated}
\begin{tabular}{|l|r|}
\hline
Quantity & Value\\
\hline
Saturated log-likelihood
    & $-1373.7676$\\
Moment upper bound $f_{\operatorname{mom},3}$
    & $-1373.7676$\\
Best log-likelihood from 100-start EM
    & $-1377.1748$\\
$f_{\operatorname{mom},3}-f_{\operatorname{EM}}$
    & $3.4072$\\
Deviance relative to the saturated model
    & $6.8144$\\
Moment relaxation runtime (seconds)
    & $39.42$\\
100-start EM runtime (seconds)
    & $29.74$\\
\hline
\end{tabular}
\end{table}
The moment-relaxation runtime covers solving the standard moment
relaxation \eqref{eq:relax_log}, whereas the EM runtime covers the
complete 100-start EM run.
The moment bound and the saturated log-likelihood agree to within
$9.92\times 10^{-8}$, which is attributable to numerical solver
tolerance. Among the 100 EM starts,
62 terminated within $10^{-6}$ of the best value, while 39 reached
the maximum of 5000 iterations. The largest shortfall from the best
EM value was $0.1843$.

This behavior follows from the geometry of the LCM image. Every point
mass on the sixteen binary outcomes belongs to the model image because
the class-conditional probabilities can be chosen deterministically.
Consequently, the convex hull of the model image is the full
probability simplex, and the convex relaxation \eqref{eq:relax_meas}
returns the saturated maximum likelihood value. Increasing the moment
relaxation order may reduce the moment-truncation gap, but it cannot
eliminate the convexification gap between this bound and the best
likelihood found over the non-saturated LCM itself.
For this instance, EM is slightly faster and returns a feasible LCM
fit, but it does not by itself certify global optimality. The moment
relaxation provides a global upper bound, but the bound is too loose to
recover or certify the non-saturated maximum likelihood estimate.
Thus, EM is more useful for obtaining an estimate in this example,
whereas the moment result serves as a diagnostic bound.
}
    
\end{example}


\section{Conclusion and Discussion}\label{sec:con}

This work explores log-polynomial optimization problems, 
characterized by logarithmic polynomial objectives with 
polynomial equality and inequality constraints. 
Using truncated $K$-moment relaxations,
we derive a hierarchy of moment relaxations
for the potentially nonconvex original formulation.
\rd{We give sufficient conditions for the relaxation to be tight and
establish criteria for detecting tightness and, under an appropriate
flat-truncation condition, extracting global optimizers.}
We further apply the Lagrange multiplier expression technique
to construct tighter relaxations 
for log-polynomial optimization with box, simplex and
ball constraints. 
\rd{The relationships among the standard moment relaxation
\eqref{eq:relax_log}, the LME moment relaxation
\eqref{eq:relax_lme}, and the original log-polynomial optimization
problem \eqref{eq:problem} are summarized in
Proposition~\ref{prop:rel}.}

The proposed approach has practical relevance in applications 
such as maximum likelihood estimation and
fixed-distribution cross-entropy or KL divergence minimization
when the predicted probabilities are polynomial in the decision variables.
Entropy-type objectives with decision-dependent coefficients
multiplying logarithms are not covered by the present framework.

\rd{Our numerical study also identifies a clear limitation of the
approach. In the tested nonconcave family, the observed difference
between the first-order standard moment upper bound
$f_{\operatorname{mom},1}$ from \eqref{eq:relax_log} and the best
feasible value
grows substantially with $m$ and varies with $n$,
whereas the size of the LME relaxation grows rapidly at similar
scales. Moreover, when a convexification gap is present,
increasing the relaxation order cannot eliminate it because the
hierarchy converges to the convex relaxation
\eqref{eq:relax_meas}.}
\rd{These findings suggest that the proposed method is most useful for
problems with small or moderate $m$ and $n$, particularly when large
or noninteger exponents make a direct full Moment--SOS relaxation
expensive or unavailable. Promising future directions include
exploiting sparsity in the moment matrices, developing less expensive
strengthenings, and extending the framework to non-polynomial
constraints.}


\section*{Statements and Declarations}

\subsection*{Funding}
\rd{Jiawang Nie is partially supported by the U.S. National Science
Foundation under grant DMS-2513254. Xindong Tang is partially supported
by the Hong Kong Research Grants Council under grants HKBU-15303423 and
HKBU-22304125.}

\subsection*{Competing interests}
\rd{The authors declare that they have no competing interests.}

\subsection*{Data availability}
\rd{The data used in the numerical experiments are either generated
by the accompanying code or included in the public GitHub repository at
\url{https://github.com/iamjiyoung/LogPOP}.}

\subsection*{Code availability}
\rd{The \textsc{Matlab} and Python code used for the numerical
experiments is available in the public GitHub repository at
\url{https://github.com/iamjiyoung/LogPOP}.}

\subsection*{Author contributions}
\rd{All authors contributed equally to this work. All authors read and
approved the final manuscript.}


\begin{thebibliography}{99}
\providecommand{\natexlab}[1]{#1}
\providecommand{\url}[1]{\texttt{#1}}
\providecommand{\urlprefix}{URL }

\bibitem{AAA13}
{\sc A.~A.~Ahmadi, A.~Olshevsky, P.~A.~Parrilo, and J.~N.~Tsitsiklis},
{\em NP-hardness of deciding convexity of quartic polynomials and related problems},
Math. Program., 137, 453--476, 2013.
\url{https://doi.org/10.1007/s10107-011-0499-2}

\bibitem{Anari24}
{\sc N.~Anari, K.~Liu, S.~O.~Gharan and C.~Vinzant},
{\em Log-concave polynomials II: High-dimensional walks and an FPRAS for counting bases of a matroid},
Ann. of Math. (2) 199 (1) 259--299, 2024.
\url{https://doi.org/10.4007/annals.2024.199.1.4}

\bibitem{Bach2025}

{\sc F.~Bach},
{\em Sum-of-squares relaxations for information theory and
variational inference},
Foundations of Computational Mathematics, 25,
pp.~865--903, 2025.
\url{https://doi.org/10.1007/s10208-024-09651-0}


\bibitem{mosek}

{\sc MOSEK ApS}, {\em MOSEK Optimization Toolbox for MATLAB Manual},
Version 10.0, 2023.
\url{https://docs.mosek.com/10.0/toolbox/index.html}


\bibitem{Banerjee2008}
{\sc O. Banerjee, L. El Ghaoui, and A. d'Aspremont},
{\em Model selection through sparse maximum likelihood estimation for multivariate Gaussian or binary data},
Journal of Machine Learning Research, 9,
pp. 485--516, 2008.
\url{https://jmlr.org/papers/v9/banerjee08a.html}


\bibitem{Bishop2006}
{\sc C. M. Bishop},
{\em Pattern recognition and machine learning. Information Science and Statistics}, Springer, New York, 2006.
\url{https://doi.org/10.1007/978-0-387-45528-0}

\bibitem{Carpenter2018}
{\sc T. Carpenter, I. Diakonikolas, A. Sidiropoulos, and A. Stewart},
{\em Near-optimal sample complexity bounds for maximum likelihood estimation of multivariate log-concave densities},
In Proceedings of the 31st Conference on Learning Theory,
Proceedings of Machine Learning Research, vol. 75,
pp. 1234--1262, 2018.
\url{https://proceedings.mlr.press/v75/carpenter18a.html}

\bibitem{Choi2024}
{\sc J. Choi, J. Nie, X. Tang, and S. Zhong},
{\em Generalized Nash equilibrium problems with quasi-linear constraints},
SIAM Journal on Optimization, 36(3), pp. 1589--1618, 2026.
\url{https://doi.org/10.1137/24M1699395}

%\bibitem{Clarke}
%{\sc F.H.~Clarke}, {\em Optimization and nonsmooth analysis},
%SIAM, Philadelphi, 1990.

%\bibitem{flatext}
%{\sc R.E. Curto and L.A. Fialkow},
%{\em Truncated k-moment problems in several variables}, Journal of Operator Theory, 54(1), pp. 189--226, 2005.

\bibitem{ChenLCM20}
{\sc H.~Chen, L.~Han, and A.~Lim},
{\em Beyond the EM algorithm: constrained optimization methods for latent class model},
Communications in Statistics--Simulation and Computation, 51,
pp.~5222--5244, 2022.
\url{https://doi.org/10.1080/03610918.2020.1764034}

\bibitem{Dempster1977}
{\sc A. P. Dempster, N. M. Laird, and D. B. Rubin},
{\em Maximum likelihood from incomplete data via the EM algorithm},
Journal of the Royal Statistical Society. Series B (Methodological),
39(1), pp. 1--22, 1977.
\url{https://doi.org/10.1111/j.2517-6161.1977.tb01600.x}

\bibitem{Fisher1922}
{\sc R. A. Fisher},
{\em On the mathematical foundations of theoretical statistics},
Philosophical Transactions of the Royal Society of London. Series A, Containing Papers of a Mathematical or Physical Character, 222, pp. 309--368, 1922.
\url{https://doi.org/10.1098/rsta.1922.0009}

\bibitem{Goodfellow2016}
{\sc I. J. Goodfellow, Y. Bengio, and A. Courville},
{\em Deep learning}, MIT Press, Cambridge, MA, 2016.
\url{https://mitpress.mit.edu/9780262035613/deep-learning/}

\bibitem{Goodfellow2014}
{\sc I. J. Goodfellow, J. Pouget-Abadie, M. Mirza, B. Xu, D. Warde-Farley, S. Ozair, A. Courville, and Y. Bengio},
{\em Generative adversarial nets},
Advances in Neural Information Processing Systems 27,
pp. 2672--2680, 2014.
\url{https://proceedings.neurips.cc/paper_files/paper/2014/hash/f033ed80deb0234979a61f95710dbe25-Abstract.html}



\bibitem{Gusmao2024}
{\sc G. S. Gusmão and A. J. Medford},
{\em Maximum-likelihood estimators in physics-informed neural networks for high-dimensional inverse problems}, Computers \& Chemical Engineering, 181, 108547, 2024.
\url{https://doi.org/10.1016/j.compchemeng.2023.108547}

\bibitem{Hansen2024}
{\sc D. Hansen, D. C. Maddix, S. Alizadeh, G. Gupta, and M. W. Mahoney},
{\em Learning physical models that can respect conservation laws},
Physica D: Nonlinear Phenomena, 457, 133952, 2024.
\url{https://doi.org/10.1016/j.physd.2023.133952}


%\bibitem{Hanson1965}
%{\sc M.A. Hanson}, {\em Inequality constrained maximum likelihood estimation}, Annals of the Institute of Statistical Mathematics, 17 (1965), pp. 311-321.

%\bibitem{HenrionKordaLasserre2020}
%{\sc D.~Henrion, M.~Korda, and J.~B. Lasserre}, {\em The Moment-SOS Hierarchy},
%  WORLD SCIENTIFIC (EUROPE), 2020.

\bibitem{GloptiPoly}
{\sc D. Henrion, J.-B. Lasserre, and J. L\"ofberg},
{\em GloptiPoly 3: Moments, optimization and semidefinite programming},
Optimization Methods and Software, 24(4--5), pp. 761--779, 2009.
\url{https://doi.org/10.1080/10556780802699201}

\bibitem{Karanam2025}
{\sc B. A. Karanam, S. Mathur, S. Sidheekh, and S. Natarajan},
{\em A unified framework for human-allied learning of probabilistic circuits}, In Proceedings of the AAAI Conference on Artificial Intelligence, 39(17), pp. 17779--17787, 2025.
\url{https://doi.org/10.1609/aaai.v39i17.33955}

\bibitem{Vinayak2019}
{\sc R. Korlakai Vinayak, W. Kong, G. Valiant, and S. M. Kakade},
{\em Maximum likelihood estimation for learning populations of parameters},
In Proceedings of the 36th International Conference on Machine Learning,
Proceedings of Machine Learning Research, vol. 97,
pp. 6448--6457, 2019.
\url{https://proceedings.mlr.press/v97/vinayak19a.html}


%\bibitem{MLEquation05}
%{\sc S. Hoşten, A. Khetan, and B. Sturmfels}, {\em Solving the Likelihood Equations}, Foundations of Computational Mathematics, 5 (2005), pp. 389--407.

%\bibitem{lasserre2015}
%{\sc J.~B. Lasserre}, {\em An Introduction to Polynomial and Semi-Algebraic
	%  Optimization}, Cambridge Texts in Applied Mathematics, Cambridge University
%  Press, 2015.

\bibitem{lasserre2001}
{\sc J.~B.~Lasserre}, {\em Global optimization with polynomials and the problem of moments}, SIAM Journal on Optimization, 11(3), pp. 796--817, 2001.
\url{https://doi.org/10.1137/S1052623400366802}

\bibitem{Lauritzen2019}
{\sc S. Lauritzen, C. Uhler, and P. Zwiernik},
{\em Maximum likelihood estimation in Gaussian models under total positivity},
The Annals of Statistics, 47(4), pp. 1835--1863, 2019.
\url{https://doi.org/10.1214/17-AOS1668}



%\bibitem{Lau09}
%{\sc M.~Laurent}, {\em Sums of squares, moment matrices and optimization over
	% polynomials}, in Emerging Applications of Algebraic Geometry, M.~Putinar and
%  S.~Sullivant, eds., The IMA volumes in mathematics and its applications, New
%  York, NY, 2009, Springer New York, pp.~157--270.

\bibitem{Yalmip}
{\sc J. Löfberg}, {\em YALMIP: A toolbox for modeling and optimization in MATLAB}, In Proceedings of the 2004 IEEE International Symposium on Computer Aided Control Systems Design, pp.~284--289, 2004.
\url{https://doi.org/10.1109/CACSD.2004.1393890}

\bibitem{Mannor2005}
{\sc S. Mannor, D. Peleg, and R. Rubinstein},
{\em The cross entropy method for classification},
In Proceedings of the 22nd International Conference on Machine Learning (ICML), pp. 561--568, 2005.
\url{https://doi.org/10.1145/1102351.1102422}

\bibitem{Mao2023}
{\sc A. Mao, M. Mohri, and Y. Zhong},
{\em Cross-entropy loss functions: theoretical analysis and applications}, In Proceedings of the 40th International Conference on Machine Learning, Proceedings of Machine Learning Research, vol.~202, pp.~23803--23828, 2023.
\url{https://proceedings.mlr.press/v202/mao23b.html}


\bibitem{McCulloch87}
{\sc C. E. McCulloch}, {\em Maximum likelihood estimation in a multinomial mixture model}, Technical Report BU-934-MA, Biometrics Unit, Cornell University, Ithaca, NY, 1987.
\url{https://ecommons.cornell.edu/bitstream/handle/1813/33048/BU-934-MA.pdf?sequence=1}

\bibitem{EMBook}

{\sc G. J. McLachlan and T. Krishnan},
{\em The EM algorithm and extensions}, 2nd ed.,
Wiley, Hoboken, NJ, 2008.
\url{https://doi.org/10.1002/9780470191613}


%\bibitem{ECM}
%{\sc X. Meng and D. B. Rubin}, {\em Maximum Likelihood Estimation via the ECM Algorithm: A General Framework}, Biometrika, 80(2) (1993), pp. 267-278.


\bibitem{MPO}
{\sc J.~Nie},
{\em Moment and polynomial optimization},
SIAM, Philadelphia, 2023.
\url{https://doi.org/10.1137/1.9781611977608}


\bibitem{MurraySAGE}
{\sc R.~Murray, V.~Chandrasekaran, and A.~Wierman},
{\em Signomial and polynomial optimization via relative entropy
and partial dualization},
Mathematical Programming Computation, 13,
pp.~257--295, 2021.
\url{https://doi.org/10.1007/s12532-020-00193-4}


%\bibitem{nie2014optimality}
%\leavevmode\vrule height 2pt depth -1.6pt width 23pt, {\em Optimality
	%  conditions and finite convergence of {L}asserre's hierarchy}, Mathematical
%  Programming, 146 (2014), pp.~97--121.



\bibitem{nie2013certifying}
{\sc J.~Nie},
{\em Certifying convergence of Lasserre’s hierarchy via flat truncation},
Mathematical Programming, 142, pp.~485--510, 2013.
\url{https://doi.org/10.1007/s10107-012-0589-9}

\bibitem{nie2014optimality}
{\sc J.~Nie}, {\em Optimality
	conditions and finite convergence of {L}asserre's hierarchy}, Mathematical
Programming, 146, pp.~97--121, 2014.
\url{https://doi.org/10.1007/s10107-013-0680-x}



\bibitem{Nie2019}
{\sc J.~Nie}, {\em Tight relaxations for polynomial optimization and {L}agrange multiplier expressions}, Mathematical Programming, 178, pp.~1--37, 2019.
\url{https://doi.org/10.1007/s10107-018-1276-2}

\bibitem{NieSunTangZ}
{\sc J. Nie, D. Sun, X. Tang, and M. Zhang},
{\em Solving polynomial variational inequality problems via Lagrange
multiplier expressions and Moment-SOS relaxations},
Computational Optimization and Applications, 90, pp.~361--394, 2025.
\url{https://doi.org/10.1007/s10589-024-00635-y}

\bibitem{nie2023convex}
{\sc J.~Nie and X.~Tang}, {\em Convex generalized {N}ash equilibrium problems
	and polynomial optimization}, Mathematical Programming, 198,
pp.~1485--1518, 2023.
\url{https://doi.org/10.1007/s10107-021-01739-7}

\bibitem{NieTangZgnep21}
{\sc J.~Nie, X.~Tang, and S.~Zhong}, {\em Rational generalized Nash
	equilibrium problems}, SIAM Journal on Optimization, 33, pp.~1587--1620, 2023.
\url{https://doi.org/10.1137/21M1456285}

\bibitem{NieWangZBilevel21}
{\sc J.~Nie, L.~Wang, J.~J.~Ye, and S.~Zhong}, {\em A Lagrange multiplier expression method for bilevel polynomial optimization}, SIAM Journal on Optimization, 31, pp.~2368--2395, 2021.
\url{https://doi.org/10.1137/20M1352375}




\bibitem{Putinar}
{\sc M.~Putinar}, {\em Positive polynomials on compact semi-algebraic sets},
Indiana University Mathematics Journal, 42, pp.~969--984, 1993.
\url{https://doi.org/10.1512/iumj.1993.42.42045}

\bibitem{Rayas2022}
{\sc A. Rayas, R. Anguluri, and G. Dasarathy},
{\em Learning the structure of large networked systems obeying conservation laws}, Advances in Neural Information Processing Systems 35, pp.~14637--14650, 2022.
\url{https://doi.org/10.52202/068431-1064}

\bibitem{Silverman1986}
{\sc B. W. Silverman},
{\em Density estimation for statistics and data analysis}, Chapman and Hall, London, 1986.
\url{https://doi.org/10.1201/9781315140919}


\bibitem{Srivastava2024}
{\sc A. Srivastava, A. Bayati, and S. M. Salapaka},
{\em Sparse linear regression with constraints: a flexible entropy-based framework}, In Proceedings of the 2024 European Control Conference (ECC), pp.~2105--2110, 2024.
\url{https://doi.org/10.23919/ECC64448.2024.10591241}

\bibitem{Zhang2018}
{\sc Z. Zhang and M. R. Sabuncu},
{\em Generalized cross entropy loss for training deep neural networks with noisy labels},
Advances in Neural Information Processing Systems 31, pp.~8778--8788, 2018.
\url{https://proceedings.neurips.cc/paper/2018/hash/f2925f97bc13ad2852a7a551802feea0-Abstract.html}



\end{thebibliography}


%%%%%%%%%%%%%%%%%
\end{document}
%%%%%%%%%%%%%%%%%
